\documentclass[11pt]{article}
\usepackage[margin=0.9in]{geometry}
\usepackage{amsmath,amssymb,amsthm,booktabs,mathrsfs}
\usepackage[T1]{fontenc}
\usepackage{lmodern,microtype}
\usepackage[colorlinks=true,linkcolor=blue,urlcolor=blue]{hyperref}
\newtheorem{theorem}{Theorem}[section]
\newtheorem{lemma}[theorem]{Lemma}
\newtheorem{proposition}[theorem]{Proposition}
\newtheorem{corollary}[theorem]{Corollary}
\theoremstyle{remark}
\newtheorem{remark}[theorem]{Remark}
\newcommand{\norm}[1]{\left\|#1\right\|}
\title{Toward a four-dimensional Yang--Mills construction\\
\large Notebook 29: total Ward extraction and classical spatial gain}
\author{Research working document}
\date{17 October 2026}
\begin{document}
\maketitle
\begin{abstract}
We continue the existence construction by treating the low field degrees
left open by the high-degree estimate of Notebook 28. The actual local
macrostep preserves output translations and the hypercubic group, under
the selected averaging map's stated permutation axiom. Its total
correction obeys the nonlinear gauge identities. Uniform kernel
locality permits a joint subsequential limit of vacuum-normalized
analytic interactions and coherent classical background families. There, the
quadratic, cubic and quartic marginal vertices are all fixed by one
coefficient multiplying a normalized gauge-invariant reference action.
The actual homogeneous classical action has the required normalization,
so this extraction has an exact algebraic coupling rebase.

We also propagate complete first-gradient H\"older regularity through
the ordinary averaging prefix, including strict support families.
Using total moment cancellations before assigning carriers, an exact
lattice factorization puts one difference on each quadratic factor.
Together with cubic and quartic spatial subtraction, this gives a
connected-activity bound of order \(C(J)L^{-\alpha}\) for the canceled
low-degree background, under the specified coherent-family and
infinite-kernel hypotheses. The coefficient \(C(J)\) is retained:
it cannot be defeated by a ratio chosen before \(J\).
Two ordinary minimum comparisons prove a classical background
semigroup modulo gauge. This allows a finite endpoint ratio
\(L=R^q\) to be chosen after fixing the elementary ratio \(R\)
and grouping \(J\).
An activity space modulo equality of normalized totals is Banach.
The full-reference extraction is bounded independently of the
small analytic radius, and classical background substitution
descends to this quotient. Combining the low- and high-degree
estimates gives a classical contraction under the stated joint
map and support hypotheses.

These results give a gauge-compatible marginal classification,
a quantitative classical localization of its remainder, and a
complete space in which the representation can change. They do
not yet establish that the actual quantum update descends to
this space with the required finite-regulator compatibility,
the required coupling trajectory,
or the remaining large-field and continuum construction.
\end{abstract}

The objective remains a complete existence proof for a nontrivial
four-dimensional Yang--Mills theory. The preceding notebooks establish
the homogeneous classical branch and a local quantum macrostep.
This notebook distinguishes exact finite-regulator local identities
from jointly extracted infinite analytic interactions and from the
operator estimates that require those coherent limiting families.
The compactness argument connects these levels at fixed parameters;
it does not provide a universal quantum update on all such limits.
Gauge invariance is required of the total functional,
not of each chosen activity coefficient.


\section{Actual macrostep symmetries and the quadratic tensor}
\label{msy-section}

This section concerns the actual local composite sector and its
common-domain analytic continuation. It does not assert equality of
different gauge sectors or spatial symmetry of an arbitrary measurable
completion outside the analytic domain. The input total density must
have the spatial symmetries under consideration.

\subsection{Symmetries of the exact local integral}
\label{msy-integral-section}

Fix the odd primitive ratio \(\ell\), a finite power \(R=\ell^s\), and
the centered nested blocks of the composite construction. Use compatible
\emph{cubic} tori, with distinct lifted interfaces as in the geometric
proof. The sufficiently large auxiliary grouping \(J\) can be chosen
as a power of \(\ell\), preserving a cofinal family of source tori.
On rectangular tori only permutations preserving the side lengths
are available.

The selected source mean is assumed to obey its permutation axiom
as well as the endpoint covariance already used in the construction.
The all-coordinate-order contour family and this axiom are the
properties stated in Ba{\l}aban, CMP 109 (1987), pp.~252--254,
especially (0.7). The implicit mean (0.10) is an admissible example;
it is not the definition of every permitted mean.

\begin{proposition}[Exact spatial symmetry of the local sector]
\label{msy-spatial-proposition}
Suppose the complete input density is invariant under the relevant
spatial transformations on the common small domain. The composite
constraint, independent direct-\(R\) internal contour gauge, physical
centered branch and actual retained link/color-ball integral commute
with
\begin{equation}
 x\longmapsto Sx+Rn,\qquad
 n\in\mathbb Z^d,\quad
 S\text{ a signed coordinate permutation}.
 \label{msy-space-group}
\end{equation}
After the output lattice is rescaled by \(R\), these give all its unit
translations and the full hypercubic point group. The same symmetries
hold for the total generated correction on the common zero-free
holomorphic domain, with its compatible logarithm branch.
\end{proposition}

\begin{proof}
Every transformation in \eqref{msy-space-group} preserves each level
of the centered hierarchy. It permutes the coordinate-order contour
family. The mean's permutation axiom therefore makes the contour
maps equivariant. Orientation reversal of an ordinary coarse bond
also respects the averaging formula. If that formula is written as
\[
 F=\exp\!\left\{N^{-1}\sum_x\log(U_xT^{-1})\right\}T,
\]
the reversed strings and reference are \(U_x^{-1},T^{-1}\), and
\begin{equation}
 \log(U_x^{-1}T)
     =-\operatorname{Ad}_{T^{-1}}\log(U_xT^{-1}),
 \qquad F_{\rm reverse}=F^{-1}.
 \label{msy-average-reversal}
\end{equation}
The primitive maps, their composition and the separate direct-\(R\)
contour map consequently commute with the transformation. The
centerline interface pivots are preserved as unoriented bonds.

The relative-coordinate calculation is essential. For
\(U_b=e^{iB_b}V_b\), a reversed bond satisfies
\begin{equation}
 U_b^{-1}
   =e^{-i\operatorname{Ad}_{V_b^{-1}}B_b}V_b^{-1}.
 \label{msy-relative-reversal}
\end{equation}
At fixed real \(V\), this is an orthogonal linear color map.
Together with the bond permutation it induces an orthogonal map
\(D_{S,V}\) on the retained coordinates. Uniqueness of the small
pivot solution gives
\begin{equation}
 \Phi_{SV}(D_{S,V}z)
       ={\cal R}_{S,V}\Phi_V(z),
 \label{msy-chart}
\end{equation}
where \({\cal R}_{S,V}\) is the full relative-bond transformation.
Thus the original window
\(\prod_{b\ {\rm retained}}1_{\{|z_b|\le gT\}}\)
is preserved with absolute Jacobian one. No complex change of a
sharp indicator is used.

The transformed physical branch is feasible, has all internal
contours equal to the identity, and has the same constrained
classical value and stationarity. Common-domain uniqueness yields
\(V_*(SW)=SV_*(W)\). The homogeneous classical action has the
required value symmetry by applying the same transformations
to its defining ordinary-constraint minimization; the fine Wilson
action and small global norm class are preserved.

Haar factors are invariant under inverses and conjugation. A pivot
derivative transforms by the relative coarse and pivot color maps.
On a reversed bond both have determinant
\((-1)^{\dim\mathfrak g}\); their adjoint factors have determinant
one. The signs cancel in the ratio. Permutations merely relabel
the factors. This accounts for both the background delta
determinant and the normalized straightening determinant.

The chart action plus absolute gauge penalty satisfies
\(T_{SV}(D_{S,V}z)=T_V(z)\). The map \(D_{S,V}\) is independent
of \(z\), so constant, linear and quadratic Taylor terms transform
separately, also in the centered independent-background extension.
The Hessian transforms by orthogonal congruence, preserving the
normalized Gaussian and its determinant. The fixed flat determinant
is the same scalar denominator. The retained change of variables
therefore proves equality of the real local integrals.

The complete inherited insertion agrees with its pre-gauge
representation on the original common domain. Its assumed
spatial invariance and holomorphic uniqueness extend the total
identity to the common-domain continuation. No symmetry of its
individual activity coefficients is required. Finally the proved
zero-free cluster domain and common branch give equality of
logarithms. Holomorphy alone would not imply this last step.
\end{proof}

Auxiliary \(J\)-grouping reorganizes the exact retained product
window and finite integral; it does not alter their scalar total.
Its fixed carrier grid and deterministic trees may break symmetry
at coefficient level. The proposition concerns their sum.
Additional genuine block-norm or asymmetric cutoff restrictions
would require a new check. A color-coordinate cube is not a
substitute for the invariant linkwise color balls.

There is an explicit obstruction to current-lattice unit
translations for the fixed sector. A pure gauge with a small
nonzero phase \(t\) at one block root and zero elsewhere has
absolute contour penalty
\[
 (R^d-1)q(e^{it}),\qquad
 q(U)=1-\tfrac12\tau(U+U^{-1}).
\]
Moving that phase one current unit to a nonroot gives only
\(q(e^{it})\). This does not conflict with output unit translations,
whose current lift is \(R\).

The measurable completion using an ordered reference contour and
fallback is gauge equivariant, but its global spatial equivariance
has not been proved. The proposition does not concern that
completion or the uncontrolled complement. If needed, averaging
transformed sector--complement decompositions preserves their
exact sum when the blocking map and input are equivariant.
Averaging conjugated asymmetric blocking kernels instead defines
a new Markov blocking kernel; it is not equality with the original
deterministic map. No such change is made here.

\paragraph{Charge conjugation.}
For \(SU(2)\), \(U\mapsto\overline U\) is conjugation by
\(i\sigma_2\), hence is already a constant gauge transformation.
It is exact on the same domains. For \(SU(N)\), the holomorphic
automorphism is \(U\mapsto(U^{-1})^T\), with orthogonal real Lie
map \(A\mapsto-A^T\). The implicit mean
\(\sum_j\log(U_jP^{-1})=0\) commutes with this automorphism
by local uniqueness. The proof then applies if the complete
input is also invariant. For a different unspecified allowed
mean, this outer-automorphism covariance is an extra property
to verify; reality alone is not a proof.

\subsection{Uniform Hessian decay and a kernel limit}
\label{msy-kernel-section}

Let the \emph{total} output correction have connected activities
on output \(J\)-cells:
\begin{equation}
 F(w)=\sum_{H\ {\rm connected}}F_H(w_H),\qquad
 \sup_I\sum_{H\ni I}e^{a_0|H|}\|F_H\|_b\le\epsilon.
 \label{msy-activity-norm}
\end{equation}
Here the supremum is on a scalar-coordinate polydisc of radius
\(b\); a color-ball domain may be replaced by a contained scalar
polydisc. Write \(D_0=d(\dim\mathfrak g)J^d\) for an upper bound
on the scalar coordinate count in a cell. All parameters are
uniform in volume and preceding homogeneous depth; \(J\) is fixed.

For bond/color coordinates \(i,j\), let \(\xi_i,\xi_j\) be their
bond midpoints and \(K_{ij}=\partial_i\partial_jF(0)\).
Cauchy gives
\[
 |\partial_i\partial_jF_H(0)|\le2b^{-2}\|F_H\|_b.
\]
If connected \(H\) contains both cells, their midpoint distance
is at most \(J(|H|+d)\), using shortest distance on a torus.
Sum the at most \(D_0|H|\) second coordinates before the
anchored support sum. With \(\mu_0=a_0/(2J)\), this proves
\begin{equation}
 \sup_i\sum_j e^{\mu_0|\xi_i-\xi_j|_1}|K_{ij}|
 \le {\cal M}:=
 \frac{2D_0e^{a_0d/2}\epsilon}{b^2}
                  \sup_{n\ge1}n e^{-a_0n/2}.
 \label{msy-exponential-row}
\end{equation}
No symmetry of the activity decomposition was used.

For the real quadratic classification assume \(F\) is real on real
small fields, as for the physical positive density and its real
logarithm. Complex inputs have the same algebraic classification
with a possibly complex \(\beta\).
For a compact semisimple group, constant gauge invariance
gives \(DF(0)=0\). Differentiate local gauge invariance once
in a site gauge and once in a field direction. The derivative
of its nonlinear generator is multiplied by \(DF(0)\), so
the exact position-space Ward identity is
\begin{equation}
 K D=0,
 \label{msy-Ward-realspace}
\end{equation}
with \(D\) the output unit gradient.

Take an exhausting sequence of compatible cubic tori and
assume the uniform bound \eqref{msy-exponential-row}.
A diagonal subsequence gives a kernel limit at every fixed
displacement. Fatou preserves exponential summability.
Translation covariance, signed permutations, Hessian symmetry
and the finite-difference identity \eqref{msy-Ward-realspace}
pass at fixed displacements. This supplies an infinite-lattice
kernel, not a probability measure, uniqueness statement or
continuum limit.

This limiting step is substantive. Finite periodic Ward
identities hold only at their allowed momenta and do not
by themselves imply a vanishing zero-momentum mass
coefficient. Small gauge-invariant winding terms are possible.
No interpolation of finite discrete momenta is used below.

\subsection{Midpoint Fourier classification with a remainder bound}
\label{msy-classification-section}

Constant adjoint invariance makes the color structure scalar
on each simple ideal. Cross terms between distinct simple
ideals vanish. Work with one such spatial kernel
\(k_{\mu\nu}(r)\). Define
\begin{equation}
 \widetilde K_{\mu\nu}(p)=
 \sum_r k_{\mu\nu}(r)
       e^{ip\cdot(r+(e_\nu-e_\mu)/2)},\qquad
 \widehat p_\mu=2\sin(p_\mu/2).
 \label{msy-midpoint}
\end{equation}
The exponential bound makes the symbol analytic near zero.
The limiting position-space Ward identity yields
\begin{equation}
 \widetilde K(p)\widehat p=0
       \quad\text{for every real }p.
 \label{msy-Ward-symbol}
\end{equation}
At \(p=te_\mu\ne0\), this kills \(e_\mu\); continuity gives
\(\widetilde K(0)=0\). Full inversion gives
\(\widetilde K(-p)=\widetilde K(p)\). Reality and Hessian
symmetry make the symbol real symmetric for real \(p\).
Midpoint coordinates give the phase-free signed-permutation
law and eliminate odd Taylor degrees.

\begin{theorem}[Quadratic Maxwell tensor]
\label{msy-Maxwell-theorem}
Under the preceding total symmetries, Ward identity and
uniform kernel hypotheses, each simple color factor has
a real coefficient \(\beta\) such that
\begin{equation}
 \widetilde K(p)
 =\beta\bigl(|p|^2I-pp^T\bigr)+O(|p|^4)
 =\beta\bigl(|\widehat p|^2I-
                      \widehat p\widehat p^T\bigr)+O(|p|^4).
 \label{msy-Maxwell-expansion}
\end{equation}
The estimate concerns the subsequential infinite kernel;
it is not an assertion about a continuous interpolation
of a finite-volume symbol.
\end{theorem}

\begin{proof}
Signed reflections force the degree-two part to have only
squares on its diagonal and \(p_\mu p_\nu\) off diagonal.
Permutations identify their coefficients:
\[
 (K_2)_{\mu\mu}=A p_\mu^2+
                  B\sum_{\nu\ne\mu}p_\nu^2,\qquad
 (K_2)_{\mu\nu}=C p_\mu p_\nu\quad(\mu\ne\nu).
\]
The cubic part of \eqref{msy-Ward-symbol} is \(K_2(p)p=0\).
Its polynomial coefficients give \(A=0\) and \(C=-B\).
Put \(\beta=B\). Analyticity and parity give the fourth-order
remainder, and \(\widehat p=p+O(|p|^3)\) gives its second form.
\end{proof}

For explicit constants set
\begin{equation}
 M_n=\max_\mu\sum_{\nu,r}
 |r+(e_\nu-e_\mu)/2|_1^n\,|k_{\mu\nu}(r)|.
 \label{msy-moments}
\end{equation}
For \(n>0\), \eqref{msy-exponential-row} gives
\[
 M_n\le{\cal M}\left(\frac{n}{e\mu_0}\right)^n,
 \qquad |\beta|\le M_2/2.
\]
The even cosine Taylor remainder yields, for real \(p\),
\[
 \|\widetilde K(p)-K_2(p)\|_{\rm op}
               \le M_4\|p\|_\infty^4/24.
\]
Using \(|2\sin(t/2)-t|\le|t|^3/24\) in every row,
for \(\|p\|_\infty\le1\) we obtain
\begin{equation}
 \left\|\widetilde K(p)
       -\beta\bigl(|\widehat p|^2I-
                        \widehat p\widehat p^T\bigr)\right\|_{\rm op}
 \le\left(\frac{M_4}{24}+\frac d6|\beta|\right)
                              \|p\|_\infty^4 .
 \label{msy-quantitative-remainder}
\end{equation}
Real symmetry turns the absolute row estimate into the
operator estimate. Without coordinate permutations there
need not be a single \(\beta\): unequal plaquette-orientation
coefficients already give a local gauge-invariant action
with translations and axis reflections but several
quadratic coefficients.

\subsection{Scope for a subsequent extraction}
\label{msy-scope-section}

The preceding argument establishes actual output-lattice
symmetry of the local \(SU(2)\) integral and the resulting
quadratic classification. It determines neither the sign
of \(\beta\) nor asymptotic freedom. The uniform flat
classical normalization is a separate, already available
input: the explicit alias formula in Notebook07 gives
\[
 q_{M,\mathrm{mid}}(p)=|p|^2I-pp^T+O(|p|^4)
\]
uniformly in \(M\). Together with the homogeneous minimizer's
identification of \(q_M\) as \(D^2S_M(0)\), this supplies
the classical quadratic normalization in that scope.
Ellipticity alone would not have supplied it.

Subtracting \(\beta\) times the normalized standard Wilson
plaquette functional is a gauge-invariant scalar subtraction
whose residual quadratic kernel satisfies
\eqref{msy-quantitative-remainder}. A localized nonlinear
extraction, the higher-field Ward completion, finite-volume
winding control and stable relevant/marginal flow remain
additional steps. Coefficientwise high-field-jet
contraction by itself does not supply them.

% No preamble. Prefix pgh-. Finite compatible regulators throughout.
\section{H\"older regularity through an ordinary prefix}
\label{pgh-section}

Fix the primitive odd ratio $\ell$ and $0<\alpha<1$. Write
$N=MR=\ell^{k+s}$, $\eta=N^{-1}$ and $h=\ell^k\eta=R^{-1}$.
Physical next-coarse spacing is one. Define the physical prefix map
\begin{equation}
 P_h(A)=\frac1{ih}\log M_\ell^k(e^{i\eta A}).
 \label{pgh-prefix}
\end{equation}
The established ordinary-trajectory theorem gives a common complex
supremum ball on which $P_h$ is holomorphic, vanishes at zero and
$\|P_h(A)\|_\infty\le K\|A\|_\infty$, uniformly in $M,R$ and
compatible periodic volume. Its dependency radius is at most
$c_\ell h$ in physical coordinates. Nested centered partitions give
\begin{equation}
 P_h(T_zA)=T_zP_h(A),\qquad z\in h\mathbb Z^d.
 \label{pgh-covariance}
\end{equation}
Only these aligned translations are asserted. Use the complete norm
\[
 \|A\|_{\mathcal N_{\alpha,\eta}}
 =\|A\|_\infty+\|\nabla_\eta A\|_\infty
                         +[\nabla_\eta A]_\alpha,
\]
with physical distances and equivalent fixed-collar versions on
periodic lifts. Constants below depend on the primitive ratio,
$\alpha$ and the group, not on $M,R$ or auxiliary grouping $J$.

\begin{lemma}[Uniform complete-field prefix bound]
\label{pgh-prefix-lemma}
On a smaller fixed ball,
\begin{equation}
 \|P_h(A)\|_{\mathcal N_{\alpha,h}}
 \le C\bigl(\|A\|_{\mathcal N_{\alpha,\eta}}
                   +\|A\|_{\mathcal N_{\alpha,\eta}}^2\bigr).
 \label{pgh-prefix-bound}
\end{equation}
The map is holomorphic into the complete output space with uniform
Banach Cauchy derivative bounds.
\end{lemma}
\begin{proof}
Uniform supremum holomorphy bounds $DP_h,D^2P_h$ on a smaller
supremum ball by constants $K_1,K_2$. For $T_h=T_{he_\mu}$,
$\nabla_hA=(T_hA-A)/h$ is an average of fine gradients; its
supremum and H\"older seminorm are bounded by those of
$\nabla_\eta A$. Equation~\eqref{pgh-covariance} gives
\begin{equation}
 \nabla_hP_h(A)=\int_0^1
 DP_h\bigl(A+t(T_hA-A)\bigr)\nabla_hA\,dt.
 \label{pgh-gradient-identity}
\end{equation}
If $\|A\|_{\mathcal N_{\alpha,\eta}}\le b$, this bounds the output
gradient by $K_1b$. Translate by $z\in h\mathbb Z^d$ and subtract.
The change in direction costs $CK_1b|z|^\alpha$. The derivative
operator changes by at most
$K_2\|T_z((1-t)A+tT_hA)-((1-t)A+tT_hA)\|_\infty$.
After multiplication by $\|\nabla_hA\|_\infty$, this costs
$CK_2b^2|z|$, hence $CK_2b^2|z|^\alpha$ when $|z|\le1$.
For larger shifts the gradient supremum suffices. This proves
\eqref{pgh-prefix-bound}, including periodic lifts and short tori
with incidences counted with multiplicity. All interpolation
segments lie in the common complex ball. At each finite regulator
the map is holomorphic between finite-dimensional normed spaces;
the displayed uniform bound on the complete-field ball and
Banach Cauchy give uniform complete-output derivative bounds.
No ordinary-average physical adjoint is used.
\end{proof}

\subsection{Supported pre-gauge fields and their differences}
\label{pgh-support-section}

Let the established ordinary representative be
$A_{\mathrm o,N}(w)=T_N(A_{\mathrm c,N}(w))$, with
$\|A_{\mathrm o,N}(w)\|_{\mathcal N_{\alpha,\eta}}
\le C_{\mathrm o}\|w\|_\infty$. Define
\begin{equation}
 a_p(w)=hP_h(A_{\mathrm o,N}(w)),\qquad
 e^{ia_p}=M_\ell^k(e^{i\eta A_{\mathrm o,N}(w)}).
 \label{pgh-pregauge}
\end{equation}
This is the selected representative before the rooted current
gauge. It is not a regularity assertion for every gauge-equivalent
representative. Lemma~\ref{pgh-prefix-lemma} implies
\begin{equation}
 \|a_p(w)\|_{\mathcal N_{\alpha,h}}
                  \le Ch\|w\|_\infty.
 \label{pgh-angle-norm}
\end{equation}
In current-unit coordinates, let $\Delta_\mu$ be the unscaled
difference between adjacent $h$-lattice sites. Thus, on $\|w\|<b$,
\begin{equation}
 \|a_p\|_\infty\le Cbh,\qquad
 \|\Delta_\mu a_p\|_\infty\le Cbh^2,\qquad
 \|\Delta_\nu\Delta_\mu a_p\|_\infty
                         \le Cbh^{2+\alpha}.
 \label{pgh-scale-bounds}
\end{equation}
More generally, for an integer current-lattice shift $z$,
\begin{equation}
 |\Delta_\mu a_p(x+z)-\Delta_\mu a_p(x)|
                  \le Cbh^{2+\alpha}|z|^\alpha.
 \label{pgh-shift-bound}
\end{equation}
For physical shifts longer than one, the gradient supremum
proves the same bound. Periodic distances use compatible lifts.
The second-difference estimate uses first-gradient H\"older
regularity; it does not assert a uniform classical $C^2$ bound.

Let support labels refer to final physical unit cells and let
$\tau$ be graph Steiner length. Fix an admissible spatial exponent
$\kappa>0$ from the classical support theorem and any prescribed
finite cardinality reserve $a\ge0$. For an output root $i$, use
one complete $h$-field fiber $E_{h,i}$ with the local
$\mathcal N_{\alpha,h}$ norm on a fixed enlarged collar of $i$,
including all required difference endpoints and H\"older quotients.

\begin{proposition}[Complete-output strict supports]
\label{pgh-support-proposition}
There are $r_a,C_a>0$, independent of $M,R$ and volume, and strict
holomorphic families
\begin{equation}
 a_p|_i=\sum_{X\ne\varnothing}a_{p;i,X}(w_X),\qquad
 a_{p;i,X}(0)=0,
 \label{pgh-family}
\end{equation}
such that, for $0<r\le r_a$,
\begin{equation}
 \sup_i\sum_X e^{\kappa\tau(\{i\}\cup X)+a|X|}
       \|a_{p;i,X}\|_{E_{h,i};r}\le C_a h r.
 \label{pgh-family-bound}
\end{equation}
The linear supported part has bound $C_ahr$ and the nonlinear
remainder has bound $C_ahr^2$, after a fixed radius reduction.
\end{proposition}
\begin{proof}
Use the classical support theorem's smooth partition and input
Banach spaces of complete supported fine fields. Reconstruction
has a uniform norm $C_{\mathrm{rec}}$; no sharply restricted
H\"older field is extended. The map $P_h\circ T_N$ is uniformly
holomorphic into $E_{h,i}$ by Lemma~\ref{pgh-prefix-lemma} and
the ordinary-transfer theorem. Its row reads at most a fixed
number $N_0$ of input physical cells in a fixed radius $r_0$.
This follows from the $O(h)$ prefix range and the fixed physical
range of $T_N$.

On a common complete-field radius $\rho_1$, Cauchy and
polarization, followed by summing only the input cell slots,
give the mixed degree-$n$ majorant
\begin{equation}
 v_n\le K
 \left(\frac{eC_{\mathrm{rec}}N_0e^{\kappa r_0}}{\rho_1}\right)^n
       =Kc^n,\qquad n\ge1.
 \label{pgh-vertex-bound}
\end{equation}
The output collar can contain $O(R^d)$ bonds, but its complete
norm is one Banach output norm, taken before the support sum.
No output coordinate count occurs in \eqref{pgh-vertex-bound}.

Substitute the corrected minimizer's complete-field support
family, whose norm with the prescribed reserve is at most
$C'_ar$. Rooted Steiner trees glue at the input slot roots,
and $e^{a|\cup_qX_q|}\le\prod_qe^{a|X_q|}$. Hence the resulting
series is bounded by $\sum_{n\ge1}K(cC'_ar)^n$, at most
$2KcC'_ar$ when $cC'_ar\le1/2$. Multiplication by $h$ proves
\eqref{pgh-family-bound}. Use one global Taylor/partition assignment
and group by the union of input variable supports before restricting
complete output coefficients to collars. Their restrictions on
overlapping collars agree; arbitrary separate decompositions at
different output roots are not used. Absolute convergence in the complete
local output norm identifies the evaluated series with
\eqref{pgh-pregauge}. Separating its linear terms gives the
quadratic remainder estimate. Applying the bounded difference
maps on the enlarged collars yields the weighted family versions
of \eqref{pgh-scale-bounds}--\eqref{pgh-shift-bound}.
\end{proof}
These are finite-regulator coefficients and can retain dependence
on the fixed regulator and baseline. Volume consistency is not
asserted. A larger prescribed reserve may shrink $r_a$.

\subsection{An oriented discrete Taylor identity}
\label{pgh-taylor-section}

For a current-unit scalar component set
$D_\nu f(v)=f(v+e_\nu)-f(v)$. Choose a fixed shortest coordinate
path $\Gamma$ from $x$ to $y$. A traversed edge has positive
direction $\nu(e)$, lower positive-coordinate endpoint $v(e)$,
and traversal sign $s(e)$. Telescoping gives
\begin{align}
 f(y)-f(x)&=\sum_\nu(y_\nu-x_\nu)D_\nu f(x)
                         +\operatorname{Rem}_{xy}(f),\notag\\
 \operatorname{Rem}_{xy}(f)
 &=\sum_{e\in\Gamma}s(e)
       [D_{\nu(e)}f(v(e))-D_{\nu(e)}f(x)].
 \label{pgh-taylor}
\end{align}
A negative traversal uses the positive difference at its lower
endpoint, not at its traversal initial point. If $n=|x-y|_1$,
\eqref{pgh-scale-bounds} and \eqref{pgh-shift-bound} imply
\begin{equation}
 |a_p(y)-a_p(x)|\le Cbh^2n,\qquad
 |\operatorname{Rem}_{xy}(a_p)|
                      \le Cbh^{2+\alpha}n^{1+\alpha}.
 \label{pgh-taylor-bound}
\end{equation}
One also has the elementary bound
$|\operatorname{Rem}_{xy}(f)|\le n^2\sup|D^2f|$.
For a winding periodic path use its lifted displacement;
no global affine function on a torus is assumed. Coherent complete
coefficients justify telescoping across collars. If a subsequent
estimate keeps a fixed starting-root spatial weight along a long
path, the exponential comparison of root weights along that path
must also be budgeted; it is not omitted by the local norm bound.

These independently derived conventions supply the scale
pattern used in CMP109, pp.~285--288,
(4.17)--(4.18) and (4.30). They do not assert the subsequent
nonlinear Ward cancellation or its local carrier assignment.

\paragraph{Constant order and scope.}
The constants above are fixed at the primitive ratio before
$R$ and $J$. Quantum coordinate-kernel estimates obtained from
activities on $J$-cells can instead contain factors depending
on their dimension and on conversion of cell distance to
current distance. Finiteness of those constants does not make
them independent of $J$. When $J$ is chosen after $R$, an
expression $C(J)R^{-\alpha}$ is therefore not yet a
noncircular contraction estimate. This regularity theorem
does not classify the low-degree kernels, supply a bounded
gauge-compatible extraction, determine a beta-function sign,
or close the quantum induction.

\section{Compactness of analytic interactions and coherent backgrounds}
\label{aic-section}

This is compactness of analytic local interactions and classical
maps at fixed parameters. It is not compactness of probability
measures or a continuum construction. It supplies a joint
infinite-volume interface for the total-kernel Ward and spatial
transport arguments.

\subsection{A limiting interaction at the original radius}
\label{aic-interaction-section}

Fix a cell size \(J\), radius \(b>0\), exponent \(a>0\), and
finitely many scalar coordinates per cell. Let compatible cubic
tori exhaust the lattice. Suppose their connected analytic
coefficients satisfy
\begin{equation}
 \sup_I\sum_{H\ni I}e^{a|H|}\|E_{N,H}\|_b\le\epsilon .
 \label{aic-input}
\end{equation}
No consistency between volumes, or symmetry of individual
coefficients, is assumed. For compactly supported fields define
the vacuum-normalized total
\begin{equation}
 \Phi_N(A)=\sum_H\bigl[E_{N,H}(A|_H)-E_{N,H}(0)\bigr].
 \label{aic-total}
\end{equation}
If \(K\) contains the cells meeting the field support, then
\(|\Phi_N(A)|\le2|K|\epsilon\) for \(\|A\|_\infty<b\).

\begin{proposition}[Interaction compactness without radius loss]
\label{aic-interaction-proposition}
There is a subsequence and a connected analytic family \(E_H\)
on the original open radius-\(b\) polydiscs with
\begin{equation}
 \sup_I\sum_{H\ni I}e^{a|H|}\|E_H\|_b\le\epsilon .
 \label{aic-limit-norm}
\end{equation}
Its vacuum-normalized total \(\Phi\) is absolutely convergent on
compact fields. On every fixed finite field support and every
radius \(r<b\), \(\Phi_N\) converges locally uniformly to \(\Phi\).
In particular all fixed derivatives of the totals converge.
\end{proposition}

\begin{proof}
Each fixed finite connected \(H\) has an unambiguous central
embedding on sufficiently large tori. Its coefficients are bounded
holomorphic functions on one finite-dimensional polydisc. Montel,
or diagonal extraction of bounded Taylor coefficients followed by
Cauchy estimates, gives a locally uniform subsequential limit.
There are countably many carriers, translates and scalar Taylor
indices, so one diagonal subsequence works for all of them.

For a finite selection of carriers, Fatou at radius \(r<b\)
passes their weighted bound. Increase \(r\) to \(b\), then
exhaust the carriers, to obtain \eqref{aic-limit-norm}.
The norm is the supremum on the open \(b\) polydisc. This
retains its original radius without asserting uniform convergence
at the boundary.

For carriers meeting \(K\), the omitted large-carrier terms obey
\begin{equation}
 \sum_{\substack{H\cap K\ne\varnothing\\|H|>Q}}
 |E_{N,H}(A|_H)-E_{N,H}(0)|
       \le2|K|\epsilon e^{-aQ}.
 \label{aic-carrier-tail}
\end{equation}
Only finitely many connected carriers of size at most \(Q\)
meet \(K\). Once the torus side is much larger than \(JQ\)
and the radius of \(K\), each has a unique central lift;
winding carriers are covered by the tail. Convergence on the
finite collection followed by \eqref{aic-carrier-tail} proves
convergence of the totals. Finite-dimensional Cauchy then
proves convergence of their fixed derivatives.
\end{proof}

The infinite kernels can therefore be derivatives of one analytic
interaction, rather than unrelated subsequential tensors.
Translations, signed permutations and reality of the totals pass
by applying the convergence to a compact field and its fixed
transform. For gauge invariance choose a compactly supported
small site gauge \(v\) such that \(A,A^v\) lie in a smaller
polydisc. The transformed field has support in the original
support plus a finite edge neighborhood. Passing the identity
\(\Phi_N(A^v)=\Phi_N(A)\) gives the limiting identity and
its differentiated Ward equations. Constant conjugation preserves
compact field support and passes as well. These are identities
of totals; no symmetry of \(E_H\) is asserted.

\subsection{Joint compactness of coherent complete-field families}
\label{aic-background-section}

Fix a preceding homogeneous depth \(M\) and endpoint ratio \(L\).
The current output spacing in final units is \(h=L^{-1}\).
Suppose the finite-regulator classical maps have the coherent
families of the supported pre-gauge construction:
\begin{align}
 a_N|_i&=\sum_{X\ne\varnothing}a_{N;i,X}(w_X),
       \qquad a_{N;i,X}(0)=0,\notag\\
 \sup_i\sum_X e^{\kappa\tau(\{i\}\cup X)+A|X|}
       \|a_{N;i,X}\|_{C^{1,\alpha};b}
       &\le Cb/L .
 \label{aic-background-norm}
\end{align}
The complete physical collar norm, including its maxima and
first-gradient H\"older quotients, is inside the support sum.
Coefficients with the same \(X\) restrict one global assignment
and agree on overlapping collars.

At fixed \(L\), each bounded collar has finitely many current
coordinates. Apply Montel and a countable diagonal extraction
to all fixed \(i,X\), output coordinates and Taylor indices.
Overlap identities pass to the limits. Each complete collar
norm is a finite collection of components and difference
quotients with fixed physical denominators. It is lower
semicontinuous under the convergence. Fatou thus preserves
\eqref{aic-background-norm} with the same constants.

The sum of terms with either
\(\tau(\{i\}\cup X)>Q\) or \(|X|>Q\) is bounded by
\[
 C(b/L)\bigl(e^{-\kappa Q}+e^{-AQ}\bigr).
\]
The remaining supports form a finite set, since they lie in
a fixed ball about \(i\). Hence the full complete fields
converge locally uniformly at each finite set of roots.
The limiting support series is normally convergent and
defines an analytic map \(a\). Uniform Cauchy remainders
also give its Banach-valued analyticity.

For compact final support \(K\), any nonzero coefficient has
\(X\cap K\ne\varnothing\). Its rooted tree therefore joins
\(i\) to \(K\), yielding
\begin{equation}
 \|a(w)\|_{\text{complete collar }i}
        \le C(b/L)e^{-\kappa\operatorname{dist}(i,K)}.
 \label{aic-background-tail}
\end{equation}
A bound summed over cells of \(K\) differs only by that finite
factor. The same estimate holds uniformly on centered finite
tori with periodic distance. It gives summable output tails
at fixed \(L\), including its fixed number of current coordinates
per unit cell.

This construction may be made jointly with the interaction
extraction and with any countable collection of fixed \(M,L\).
A cofinal sequence of tori can be chosen divisible by the first
finitely many required powers at each diagonal stage. The
result supplies coherent infinite families; it does not choose
a unique background independently of the subsequence.

\subsection{Summable fields and actual scalar composition}
\label{aic-composition-section}

The limiting interaction extends to summable coordinate fields
with \(\|A\|_\infty\le\rho<b\). Coordinate Cauchy and the
segment \(tA\) give
\[
 |E_H(A|_H)-E_H(0)|
 \le\frac{C\|E_H\|_b}{b-\rho}
                         \sum_{j\in H}|A_j|.
\]
Summing first over carriers containing each coordinate yields
\begin{equation}
 |\Phi(A)|\le\frac{C\epsilon}{b-\rho}\|A\|_{\ell^1},
 \qquad
 |\Phi(A)-\Phi(B)|
       \le\frac{C\epsilon}{b-\rho}\|A-B\|_{\ell^1}
 \label{aic-interaction-Lipschitz}
\end{equation}
when both fields lie in that interior ball. Coordinate counts
are already in the \(\ell^1\) sum; there is no volume factor.
The same estimates hold on the finite regulators.

Choose \(L\) so the classical amplitude bound places \(a_N(w)\)
in a fixed interior radius \(\rho<b\) on the final data domain.
For compact \(w\), \eqref{aic-background-tail} gives uniform
\(\ell^1\) tails of \(a_N(w)\). On centered representatives,
the periodic distance from an output root to \(K\) is at least
its centered distance from the origin minus the fixed radius
of \(K\), so this tail bound is uniform in volume.
Truncate the output to a large fixed region. On the region,
both the interaction coefficients and the classical fields
converge locally uniformly. Outside it use
\eqref{aic-interaction-Lipschitz}. It follows that
\begin{equation}
 \Phi_N(a_N(w))\longrightarrow\Phi(a(w))
 \label{aic-substitution-limit}
\end{equation}
locally uniformly on compact final supports and smaller data
radii. This is convergence of actual scalar compositions,
not only of formal Taylor series.

For decaying noncompact intermediate inputs, apply coordinate
Cauchy directly to \eqref{aic-background-norm}. Only supports
containing the input coordinate \(j\) contribute, and their
rooted tree spans \(i,j\). On an interior ball,
\begin{equation}
 \|\partial_{w_j}a_i\|
     \le\frac{C}{b-\rho}
                   e^{-\kappa\operatorname{dist}(i,j)}.
 \label{aic-map-derivative}
\end{equation}
The fixed input and output fiber dimensions, including \(L^4\),
are absorbed in \(C\). Summing over output roots proves an
\(\ell^1\)-to-\(\ell^1\) derivative bound, uniform in volume
at fixed \(L\). Segment integration gives the corresponding
Lipschitz estimate. Since \(a(0)=0\), summable inputs have
summable outputs.

The same truncation proof therefore applies to finite
compositions, provided every intermediate field remains in
a strict common interior of its analytic domain. Finite-volume
scalar identities, including a classical semigroup identity,
can pass to a jointly extracted finite list of maps and
interactions on compatible nested volumes. Separate arbitrary
subsequences need not compose.

\subsection{Scope}
\label{aic-scope}

The result jointly attains a subsequential analytic interaction
and coherent classical families at fixed parameters. It supplies
the total-kernel and complete-field hypotheses used by the
spatial remainder theorem. It proves neither uniqueness,
convergence of Gibbs measures, a continuum theory, nor consistency
of a quantum update for all input families. Descent of the actual
quantum map to a chosen interaction space, continuity there with
its normalizations and cutoffs retained, and infinite iteration
remain separate requirements.

\section{Total Ward extraction through field degree four}
\label{tkw-section}

The following argument completes the low-degree \emph{total-kernel}
classification. It does not impose gauge invariance on individual
activity coefficients. Its exact statement concerns subsequential
infinite-lattice kernels; finite-period winding terms and localization
of the canceled functional remain separate questions.

\subsection{Hypotheses and the reference action}
\label{tkw-hypotheses}

For a compact connected simple matrix group, use Hermitian Lie-algebra
coordinates and a fixed invariant positive trace form. Consider analytic
total corrections \(E_N\) on exhausting compatible cubic tori. Assume
the local gauge identity, output-unit translations, full signed
permutations and reality on real fields established for the actual
local sector in Section~\ref{msy-section}. Assume connected activity
representations on a fixed cell partition with
\begin{equation}
 \sup_I\sum_{H\ni I}e^{a_0|H|}\|E_{N,H}\|_b\le\epsilon .
 \label{tkw-activity}
\end{equation}
Here \(a_0,b>0\), the cell size \(J\), and the scalar coordinate count
\(D\) per cell are fixed uniformly in volume and preceding refinement.
The total symmetries need not hold for the chosen coefficients.

Let \(S_N\) be a real analytic, exactly gauge-invariant and spatially
invariant reference, with \(S_N(0)=DS_N(0)=0\), and with uniform
fixed-order kernel moment bounds through degree four. Assume its
limiting midpoint Hessian has the normalization
\begin{equation}
 \Gamma_2^S(p)=|p|^2I-pp^T+O(|p|^4).
 \label{tkw-reference}
\end{equation}
The normalized Wilson action is one choice. The actual homogeneous
classical action \(S_{MR}\) is another: its Hessian is \(q_{MR}\),
whose uniform normalization was proved by the alias formula in
Notebook07, and the classical strict-support construction supplies
the higher fixed-order bounds. Ellipticity alone would not establish
\eqref{tkw-reference}.

\begin{theorem}[Full marginal completion from the total Ward identity]
\label{tkw-theorem}
Under these assumptions, choose a joint subsequential infinite-lattice
limit of the kernels of \(E_N,S_N\) through degree four. There is a real
coefficient \(\beta\) such that the midpoint Fourier vertices of the
residual \(F=E-\beta S\) satisfy
\begin{align}
 \Gamma_2^F(p)&=O(|p|^4),\notag\\
 \Gamma_3^F(p_1,p_2,p_3)
     &=O\!\left(\max_j|p_j|^3\right),\notag\\
 \Gamma_4^F(p_1,p_2,p_3,p_4)
     &=O\!\left(\max_j|p_j|^2\right),
 \qquad \sum_jp_j=0.
 \label{tkw-remainders}
\end{align}
Thus the total Taylor coefficients of field degree \(n=2,3,4\)
and momentum degree at most \(4-n\) agree with \(\beta S\).
The one-point kernel vanishes; the vacuum is unrestricted.
\end{theorem}

For a semisimple product there is one coefficient per simple ideal.
The analogous statement subtracts gauge-invariant references with the
individual normalized Maxwell blocks. A single-coupling rebase is
immediate for a simple group, or if those coefficients coincide.
Reality can be omitted if a complex coefficient \(\beta\) is permitted.

\subsection{Fixed-order kernels and their joint limit}
\label{tkw-kernel-limit}

Write \(K_n=D^nE_N(0)\). Coordinate Cauchy on one coefficient gives
\begin{equation}
 \sum_{i_1,\ldots,i_n\in H}
  |D^nE_{N,H}(0)[e_{i_1},\ldots,e_{i_n}]|
 \le (D|H|+n)^n b^{-n}\|E_{N,H}\|_b .
 \label{tkw-coordinate-Cauchy}
\end{equation}
Indeed the exact monomial count multiplied by \(n!\) is
\(n!\binom{D|H|+n-1}{n}\), bounded by the displayed factor.
Fixing the first coordinate only decreases the sum. The graph
Steiner length of slots contained in connected \(H\) is at most
\(|H|-1\). For \(0<\sigma<a_0\), therefore,
\begin{equation}
 \sup_{i_1}\sum_{i_2,\ldots,i_n}
 e^{\sigma\tau(i_1,\ldots,i_n)}
 |K_n(i_1,\ldots,i_n)|
 \le b^{-n}\epsilon
       \sup_{m\ge1}(Dm+n)^n e^{-(a_0-\sigma)m}.
 \label{tkw-kernel-bound}
\end{equation}
Scalar and color indices are included in the sums. A further fixed
exponent loss absorbs every fixed spatial moment. Conversion from
the \(J\)-cell metric to the unit-site metric retains \(J\)-dependent
constants.

Global constant gauges imply \(K_1=0\) separately at each bond,
because its covector annihilates every Lie commutator. Translation
invariance permits an anchor at the origin. The bounds then give a
diagonal subsequence of all fixed displacements, simultaneously for
\(n\le4\) and for the reference. Fatou preserves exponential
summability. Permutation symmetry of the slots, signed permutations,
translation and global adjoint covariance pass to this limit.
The differentiated Ward identities below pass as well: their
coefficient equations contain finite nearest-neighbor shifts and
lower kernels with coinciding slots. Equivalently one may first
test against finitely supported fields and gauges.

The limiting Fourier vertices are analytic near the simultaneous
origin, and their fixed derivatives are justified by the spatial
moments. No unnormalized infinite-volume total functional is being
summed; even a nonzero vacuum density would preclude that.

\subsection{The exact logarithmic gauge generator}
\label{tkw-generator-section}

For \(U_b=e^{iA_b}\), use
\[
 U_b(t)=e^{it\lambda_{b_-}}U_be^{-it\lambda_{b_+}},
 \qquad (d\lambda)_b=\lambda_{b_+}-\lambda_{b_-}.
\]
Differentiating the exponential chart, with
\(\operatorname{ad}_A X=[A,X]\), gives
\begin{equation}
 \delta_\lambda A
 =i[\lambda_-,A]-f(i\operatorname{ad}_A)d\lambda,
 \qquad
 f(z)=\frac{z}{1-e^{-z}}
       =1+\frac z2+\frac{z^2}{12}+O(z^4).
 \label{tkw-generator}
\end{equation}
In particular
\begin{align}
 \delta_\lambda A
   &=-d\lambda+L_\lambda(A)+Q_\lambda(A,A)
                        +O(A^4\lambda),\notag\\
 L_\lambda(A)_b
   &=\frac i2[\lambda_{b_-}+\lambda_{b_+},A_b],\notag\\
 Q_\lambda(A,B)_b
   &=\frac1{24}\bigl(
      [A_b,[B_b,d\lambda_b]]+
      [B_b,[A_b,d\lambda_b]]\bigr).
 \label{tkw-LQ}
\end{align}
The cubic Bernoulli coefficient is zero. Comparing powers in
\(DE(A)[\delta_\lambda A]=0\), using \(K_1=0\), gives
\begin{align}
 K_2(A,d\lambda)&=0,\notag\\
 K_3(A,A,d\lambda)&=2K_2(A,L_\lambda(A)),\notag\\
 K_4(A,A,A,d\lambda)
   &=3K_3(A,A,L_\lambda(A))
       +6K_2(A,Q_\lambda(A,A)).
 \label{tkw-Ward-diagonal}
\end{align}
Polarization yields
\begin{align}
 K_3(A,B,d\lambda)
   &=K_2(A,L_\lambda(B))+K_2(B,L_\lambda(A)),\label{tkw-Ward-cubic}\\
 K_4(A,B,C,d\lambda)
   &=K_3(A,B,L_\lambda(C))
      +K_3(A,C,L_\lambda(B))
      +K_3(B,C,L_\lambda(A))\notag\\
   &\quad+2K_2(A,Q_\lambda(B,C))
      +2K_2(B,Q_\lambda(A,C))
      +2K_2(C,Q_\lambda(A,B)).
 \label{tkw-Ward-quartic}
\end{align}
These are total-kernel identities. Subtracting just a quadratic
polynomial would generally violate them. Subtracting a multiple
of an exactly invariant reference preserves them.

The generator and differentiated identities are independently
derived here. Their published counterparts are Ba{\l}aban,
CMP 109 (1987), (4.9)--(4.15), pp.~283--284; the full-action
subtraction is (4.38)--(4.45), pp.~291--292. The stronger
coefficientwise orbit hypotheses of that article are not assumed.

\subsection{Midpoint phases and parity}
\label{tkw-midpoint-section}

Use bond-midpoint Fourier fields and site Fourier gauges:
\[
 A_\mu(x)=\int e^{ip\cdot(x+e_\mu/2)}a_\mu(p)\,dp,
 \qquad \lambda(x)=\int e^{ik\cdot x}\lambda(k)\,dk .
\]
On \(\sum_jp_j=0\), anchoring the first bond gives the vertex phase
\begin{equation}
 \exp\!\left\{i\sum_{j=2}^n
   p_j\cdot(x_j+e_{\mu_j}/2-e_{\mu_1}/2)\right\}.
 \label{tkw-vertex-phase}
\end{equation}
Only local analyticity near zero momentum is required; midpoint
vertices need not be individually periodic under every \(2\pi\)
shift. The Fourier coefficients in \eqref{tkw-LQ} are exactly
\begin{align}
 (d\lambda)_\mu&:\quad
       2i\sin(k_\mu/2)\lambda(k),\notag\\
 L_\lambda(A)_\mu&:\quad
       i\cos(k_\mu/2)[\lambda(k),a_\mu(p)],\notag\\
 Q_\lambda(A,B)_\mu&:\quad
       \frac{i}{12}\sin(k_\mu/2)
       \bigl([a_\mu(p),[b_\mu(q),\lambda(k)]]
           +[b_\mu(q),[a_\mu(p),\lambda(k)]]\bigr).
 \label{tkw-generator-phases}
\end{align}
Their output momenta are \(k,p+k,p+q+k\). Thus \(L\) has
a bounded analytic coefficient, while \(Q\) supplies one
power of gauge momentum.

Full inversion sends \(a_\mu(p)\) to \(-a_\mu(-p)\).
Consequently
\begin{equation}
 \Gamma_n(-p_1,\ldots,-p_n)
        =(-1)^n\Gamma_n(p_1,\ldots,p_n).
 \label{tkw-parity}
\end{equation}
The cubic vertex is odd; the quadratic and quartic vertices
are even. Midpoint phases are essential for this form.

\subsection{Proof of the marginal completion}
\label{tkw-completion-proof}

\begin{proof}[Proof of Theorem~\ref{tkw-theorem}]
Theorem~\ref{msy-Maxwell-theorem} gives
\[
 \Gamma_2^E(p)=\beta(|p|^2I-pp^T)+O(|p|^4).
\]
For completeness, the continuous limiting Ward relation
\(\Gamma_2(p)\widehat p=0\) first implies \(\Gamma_2(0)=0\)
by taking \(p=te_\mu\). Reflections and permutations restrict
the degree-two tensor to diagonal
\(A p_\mu^2+B\sum_{\nu\ne\mu}p_\nu^2\) and off-diagonal
\(C p_\mu p_\nu\). Transversality gives \(A=0,C=-B\).
Put \(\beta=B\), which is real under the stated reality
assumption. It can be extracted as
\begin{equation}
 \beta=\frac12\,
  \partial_{p_\nu}^2(\Gamma_2^E)_{\mu\mu}(0),
       \qquad \mu\ne\nu,
 \label{tkw-beta}
\end{equation}
in a trace-orthonormal color direction. The reference
normalization proves \(\Gamma_2^F(p)=O(|p|^4)\).

Apply \eqref{tkw-Ward-cubic} to the residual and momenta
\(tp,tq,tk\), with \(p+q+k=0\). Its right side is \(O(t^4)\):
each residual Hessian is evaluated at small momentum, and
the \(L\) coefficient in \eqref{tkw-generator-phases} is bounded.
Parity gives
\[
 \Gamma_3^F(tp,tq,tk)=tP_1(p,q,k)+O(t^3).
\]
The left side contracts its last vector index with
\(2i\sin(tk/2)=itk+O(t^3)\). Its degree-two coefficient
therefore vanishes. Slot permutation symmetry gives, for
each choice of colors,
\begin{equation}
 p_\mu(P_1)_{\mu\nu\rho}=0,\qquad
 q_\nu(P_1)_{\mu\nu\rho}=0,\qquad
 k_\rho(P_1)_{\mu\nu\rho}=0.
 \label{tkw-linear-transverse}
\end{equation}

Eliminate \(k=-p-q\) and write
\[
 (P_1)_{\mu\nu\rho}
  =\sum_\sigma A_{\mu\nu\rho;\sigma}p_\sigma
                   +B_{\mu\nu\rho;\sigma}q_\sigma .
\]
In the first relation of \eqref{tkw-linear-transverse},
the independent mixed monomials \(p_\mu q_\sigma\)
force every \(B_{\mu\nu\rho;\sigma}\) to vanish.
The second relation then forces every \(A_{\mu\nu\rho;\sigma}\)
to vanish by the independent monomials \(q_\nu p_\sigma\).
Thus \(P_1=0\), proving the cubic remainder. This elementary
argument permits arbitrary complex/color coefficients and
does not presume a separately invariant cubic functional.

Finally apply \eqref{tkw-Ward-quartic}, scaling all four
momenta by \(t\). Its \(K_3,L\) terms are \(O(t^3)\);
its \(K_2,Q\) terms are \(O(t^5)\), because \(Q\) has
the sine factor. By parity write
\[
 \Gamma_4^F(tp_1,\ldots,tp_4)=P_0+O(t^2).
\]
The order-\(t\) term on the left contracts \(P_0\) with
the last momentum \(p_4\). That momentum is arbitrary:
the other three can always sum to \(-p_4\). Hence
every component of \(P_0\) vanishes. This proves the
quartic remainder and completes the theorem.
\end{proof}

The argument is an identity of the integrated vertices.
It does not choose a unique density modulo lattice divergences.

\subsection{Reference independence and the full action}
\label{tkw-full-action-section}

For the standard Wilson reference, the marginal continuum
polynomial is explicit. With Hermitian coordinates set
\[
 f_{\mu\nu}=\partial_\mu A_\nu-\partial_\nu A_\mu,
 \qquad F_{\mu\nu}=f_{\mu\nu}+i[A_\mu,A_\nu].
\]
Its engineering-degree-four part is
\begin{equation}
 \frac12\sum_{\mu<\nu}\tau(F_{\mu\nu}^2)
 =\frac12\sum_{\mu<\nu}\tau(f_{\mu\nu}^2)
   +i\sum_{\mu<\nu}\tau(f_{\mu\nu}[A_\mu,A_\nu])
   -\frac12\sum_{\mu<\nu}\tau([A_\mu,A_\nu]^2).
 \label{tkw-Wilson-polynomial}
\end{equation}
Indeed bond angles \(tA\) evaluated at midpoints of a
lattice of spacing \(t\) give the plaquette
\(P=I+it^2F_{\mu\nu}+O(t^3)\). Therefore
\(1-\frac12\tau(P+P^{-1})
 =\frac12t^4\tau(F_{\mu\nu}^2)+O(t^5)\).
This verifies the normalization and the relative cubic
and quartic coefficients.

No explicit Wilson expansion is needed in the theorem.
Applying it to the difference of two exactly invariant
normalized references shows that their marginal total
vertices agree. In particular one may subtract the actual
homogeneous classical reference directly:
\begin{equation}
 e^{-g^{-2}S_{MR}+E}
  =e^{-(g^{-2}-\beta)S_{MR}+(E-\beta S_{MR})}.
 \label{tkw-coupling-rebase}
\end{equation}
The vacuum is separate. This identity neither proves
positivity of \(g^{-2}-\beta\) nor the required sign,
size or accumulated drift of its actual increments.

\subsection{Moment remainders and the remaining norm question}
\label{tkw-scope-section}

For the residual anchored kernels use the \emph{bond midpoints}
\(\xi_j=x_j+e_{\mu_j}/2\), and define
\begin{equation}
 T_{n,m}=\sup_{i_1}\sum_{i_2,\ldots,i_n}
  \left(\sum_{j=2}^n|\xi_j-\xi_1|_1\right)^m
                  |K_n^F(i_1,\ldots,i_n)|.
 \label{tkw-moments}
\end{equation}
They are bounded by the moments of \(E\) plus
\(|\beta|\) times those of \(S\). The phase
\eqref{tkw-vertex-phase} and the real one-variable
Taylor integral remainder give componentwise
\begin{align}
 |\Gamma_2^F(p)|
   &\le \frac{T_{2,4}}{4!}\|p\|_\infty^4,\notag\\
 |\Gamma_3^F(p_1,p_2,p_3)|
   &\le \frac{T_{3,3}}{3!}
                  \max_{j\ge2}\|p_j\|_\infty^3,\notag\\
 |\Gamma_4^F(p_1,p_2,p_3,p_4)|
   &\le \frac{T_{4,2}}{2!}
                  \max_{j\ge2}\|p_j\|_\infty^2 .
 \label{tkw-quantitative}
\end{align}
Fixed vector/color row sums add only their fixed dimensions.
The exponential kernel bound makes these moments finite.

The cancellations also hold in ordinary tail-anchored
coordinates: conversion multiplies each midpoint component
by a nonzero analytic phase and cannot lower its vanishing
order. Thus all residual ordinary momentum jets through
degrees \(3,2,1\) vanish for field degrees \(2,3,4\),
respectively. This includes the weaker anchored moments
used in second-order quadratic, first-order cubic and
zeroth-order quartic spatial subtractions.

The constants can depend strongly on the fixed old grouping
\(J\). Since the common-domain construction chooses \(J\)
after \(R\), an estimate \(C(J)R^{-\alpha}\) is not
automatically a contraction. No such claim is made.
The remaining tasks are localization of the canceled
total without false coefficientwise Ward assumptions,
finite-period wrapping control or a consistent infinite
correction, a spatial-gain norm with the correct parameter
order, and the actual relevant/marginal recursion.
The theorem proves the total-kernel Taylor classification;
it does not supply these separate conclusions.

% No preamble. Prefix csg-. Scalar classical background identities only.
\section{Classical background composition modulo gauge}
\label{csg-section}

Fix the primitive odd ratio $\ell$ and write
$M=\ell^k$, $R_1=\ell^{s_1}$ and $R_2=\ell^{s_2}$.
Let $U_n(D)$ be the selected ordinary homogeneous Wilson
minimizer at refinement $n$. The ordinary-minimum theorem
of Notebook15 gives a refinement-uniform small global
$C^{1,\alpha}$ comparison ball, a bound
$\|A_{\mathrm o,n}(D)\|_{C^{1,\alpha}}
\le C\|\log D/i\|_\infty$, and gauge equivalence whenever
equality holds in that comparison class.

Define the pre-root-gauge map and its scalar substitution by
\begin{equation}
 P_{M,R}(D)=M_\ell^k(U_{MR}(D)),\qquad
 (\mathcal T_{M,R}E)(D)=E(P_{M,R}(D)).
 \label{csg-maps}
\end{equation}
Here $R=\ell^s$ and $M_\ell^sP_{M,R}(D)=D$.
All expressions include the natural identification of the
nested lattices. The input $D$ lives on the next-coarse
lattice and $P_{M,R}(D)$ on its $R$-times finer current
lattice. The prefix estimate gives a common classical
radius $b_{\mathrm{cl}}$ and constant $C_0$, independent of
$M,R$ and compatible volume, such that
\begin{equation}
 \left\|\frac1i\log P_{M,R}(e^{iw})\right\|_\infty
                \le\frac{C_0}{R}\|w\|_\infty.
 \label{csg-amplitude}
\end{equation}
The operator $\mathcal T$ is a classical background
substitution, not a fluctuation integral.

\subsection{Two minimum comparisons on the same fine graph}
\label{csg-comparison-section}

For a final datum $Z$ put
\begin{equation}
 U_2=U_{MR_1R_2}(Z),\qquad
 D_1=M_\ell^{k+s_1}(U_2)=P_{MR_1,R_2}(Z),\qquad
 U_1=U_{MR_1}(D_1).
 \label{csg-fields}
\end{equation}
Both fine fields live on the same original oriented graph.
In final-$Z$ units its spacing is
$\eta_2=(MR_1R_2)^{-1}$ and the $D_1$ spacing is $R_2^{-1}$.
The definition of $U_{MR_1}(D_1)$ instead uses $D_1$-unit
coordinates and spacing $\eta_1=(MR_1)^{-1}=R_2\eta_2$.
This is a coordinate relabeling of the same links. In four
dimensions the coupling-free Wilson action is the same
plaquette sum in both conventions.

Exact ordinary composition gives
\begin{align}
 M_\ell^{k+s_1}(U_1)&=D_1=M_\ell^{k+s_1}(U_2),\notag\\
 M_\ell^{k+s_1+s_2}(U_1)&=Z=M_\ell^{k+s_1+s_2}(U_2).
 \label{csg-feasibility}
\end{align}
For fixed finite $R_1,R_2$ both fields lie in the needed
comparison balls after shrinking the real $Z$ radius,
uniformly in $M$ and volume. Explicitly, a fine physical
potential $A_1$ in $D_1$ units becomes
\begin{equation}
 A_1^Z(x)=R_2A_1(R_2x),\qquad
 \|A_1^Z\|_{C^{1,\alpha}}
          \le C R_2^{2+\alpha}\|A_1\|_{C^{1,\alpha}}.
 \label{csg-rescaling}
\end{equation}
The reverse conversion decreases the norm up to a fixed
convention constant, and decreases its supremum by $R_2^{-1}$.
By \eqref{csg-amplitude}, $\|\log D_1/i\|\le C_0\|\log Z/i\|/R_2$.
Thus $U_1$ in final units is bounded by
$CR_2^{1+\alpha}\|\log Z/i\|$, while $U_2$ in $D_1$ units
is bounded by $C\|\log Z/i\|/R_2$. These are finitely many
radius restrictions for fixed ratios.

The minimum at $D_1$, with competitor $U_2$, and the minimum
at $Z$, with competitor $U_1$, yield respectively
\begin{equation}
 S(U_1)\le S(U_2),\qquad S(U_2)\le S(U_1).
 \label{csg-comparisons}
\end{equation}
Equality and the ordinary-minimum equality clause give a
fine gauge $g$ with $U_1=U_2^g$. This gauge is supplied by
the small ordinary-to-corrected transfer, the
constraint-preserving flow into the flat slice, and the
inverse transfer from the same strict slice minimizer.
These controlled gauge transformations can be joined to
the identity inside a smaller common neighborhood.
Equal action alone, without that equality clause, would
not justify this conclusion.

Ordinary gauge covariance gives
\begin{equation}
 M_\ell^k(U_1)=\bigl(M_\ell^k(U_2)\bigr)^{g|_{\mathrm{current}}}.
 \label{csg-gauge-covariance}
\end{equation}
Consequently
\begin{equation}
 P_{M,R_1}(P_{MR_1,R_2}(Z))
       \quad\hbox{is gauge equivalent to}\quad
 P_{M,R_1R_2}(Z).
 \label{csg-orbit-composition}
\end{equation}
The two selected representatives need not coincide. The
connecting gauge need not be unique; its restriction to
$D_1$ sites stabilizes $D_1$.

\subsection{The scalar germ identity and its continuation}
\label{csg-identity-section}

\begin{theorem}[Classical background semigroup modulo gauge]
\label{csg-theorem}
Let $E$ be a scalar total functional holomorphic on a
current-angle product ball and locally gauge invariant
there. On finite compatible regulators,
\begin{equation}
 \mathcal T_{MR_1,R_2}(\mathcal T_{M,R_1}E)
                   =\mathcal T_{M,R_1R_2}E
 \label{csg-semigroup}
\end{equation}
as holomorphic germs at zero. The identity extends to every
connected common product domain of the specified branches.
The output $\mathcal T_{M,R}E$ is locally gauge invariant.
In particular, for any fixed finite $q$,
\begin{equation}
 \mathcal T_{MR^{q-1},R}\cdots
 \mathcal T_{MR,R}\mathcal T_{M,R}E
                         =\mathcal T_{M,R^q}E.
 \label{csg-iteration}
\end{equation}
If $R>C_0$ and $b$ lies in the common classical and original
$E$ domains, both sides of \eqref{csg-iteration} are
analytic on the same radius-$b$ ball and agree there.
\end{theorem}
\begin{proof}
On a sufficiently small real $Z$ ball, the gauge path in
\eqref{csg-gauge-covariance} remains inside the original
domain of $E$. Total gauge invariance proves
\eqref{csg-semigroup}. Both compositions are holomorphic;
uniqueness from the real slice proves the germ identity
and its continuation. No gauge identity for separate
activity coefficients is used.

To see local invariance of the output, extend a small
coarse gauge smoothly to prescribed fine point values.
The transformed minimum at $D$ is a small feasible
competitor at $D^v$. Compare with the minimum at $D^v$,
and perform the inverse-gauge comparison at $D$.
Equality follows; the ordinary-minimum equality clause
identifies the two configurations modulo fine gauge.
Their partial averages are gauge equivalent, so evaluating
$E$ gives local invariance. Analyticity continues this
identity from the real overlap.

The identity can now be iterated, proving
\eqref{csg-iteration}. If each ratio is larger than $C_0$,
\eqref{csg-amplitude} maps the radius-$b$ ball into itself
at every substitution; the direct map has the same
property. Thus both sides are analytic on the connected
common $b$ ball. Their germ equality continues there,
even if the real variational comparison required a
smaller radius depending on the finite ratios or on $q$.
This uses no complex minimum assertion.
\end{proof}

\subsection{Use and limitations of a finite iterate}
\label{csg-scope-section}

The direct representative in \eqref{csg-iteration} obeys
the H\"older prefix estimates
\begin{equation}
 \|a_{M,R^q}\|_\infty\le CbR^{-q},\qquad
 \|\Delta a_{M,R^q}\|_\infty\le CbR^{-2q},\qquad
 \|\Delta^2a_{M,R^q}\|_\infty
                         \le CbR^{-(2+\alpha)q},
 \label{csg-iterated-scales}
\end{equation}
with primitive constants independent of $q$. Its complete
supported families follow by applying the same prefix
theorem at endpoint ratio $R^q$.

This permits keeping the elementary $R$ and grouping $J$
fixed before selecting a finite $q$. A separately proved
bound $C(J)R^{-q\alpha}$ could then have a noncircular
parameter order. Such an activity-norm estimate is not
proved by the scalar identity alone.

The canonical families from successive core/tail
transports need not equal the direct family. Their scalar
sums agree; a new justified localization may represent
that common scalar function. No automatic bound for the
old sequential family, or bounded relocalization on a
chosen quotient norm, follows.

Nor does \eqref{csg-iteration} identify sequential quantum
integrals with one integral at ratio $R^q$. Sequential
steps retain their separate gauge factors, normalizations,
windows and newly generated interactions. Those terms do
not disappear from an iterated quantum recursion. This
is a finite classical background semigroup modulo gauge,
uniform in preceding refinement and compatible volume;
infinite iteration and regulator consistency remain open.

\section{Connected localization of spatially canceled low kernels}
\label{lst-section}

This section proves a connected-activity norm estimate for the
canceled \emph{total} kernels. Its inputs are coherent complete-field
support families and exact infinite-lattice moment cancellations.
It also applies to explicitly periodized versions of those inputs,
with their lifted sums. It does not silently impose lifted moment
identities on an arbitrary finite torus kernel.

\subsection{Inputs and parameter order}
\label{lst-input-section}

Work in dimension four. A current coordinate consists of a site,
positive bond direction and color, with fixed internal dimension \(v\).
For \(n=2,3,4\), let \(K_n\) be a translation-invariant symmetric
\(n\)-linear kernel. Anchor its first site at \(x\), write the other
sites as \(x+r_2,\ldots,x+r_n\), and put
\(m=\sum_{j=2}^n|r_j|_1\). Internal indices are included in all sums.
Assume
\begin{equation}
 \sup_{\text{first internal index}}
 \sum_{\text{other indices},r_2,\ldots,r_n}
       e^{\sigma m}|K_n(0,r_2,\ldots,r_n)|
       \le \epsilon b^{-n}.
 \label{lst-kernel-bound}
\end{equation}
The constants \(\sigma,\epsilon\) can incorporate the fixed old
grouping \(J\). Fixed-order Cauchy and a connected activity bound
supply this estimate with a smaller spatial exponent.

Use tail-site coordinates. Assume componentwise
\begin{align}
 \sum_r r^\gamma K_2(r)&=0
                          &&(|\gamma|\le3),\notag\\
 \sum_{r,s}K_3(r,s)&=0,\qquad
 \sum_{r,s}r_\rho K_3(r,s)=
 \sum_{r,s}s_\rho K_3(r,s)=0
                          &&(\rho=1,\ldots,4),\notag\\
 \sum_{r,s,t}K_4(r,s,t)&=0.
 \label{lst-moments}
\end{align}
These are consequences of the stronger vanishing orders in
Theorem~\ref{tkw-theorem}; analytic conversion from midpoint
to tail coordinates preserves those orders.

Fix an integer endpoint ratio \(L\), put \(h=L^{-1}\), and let
\(a(w)\) be an analytic current-angle background on the radius-\(b\)
final-field ball. Current sites have physical spacing \(h\).
Assume coherent complete-field strict families: on every fixed
enlarged final-unit collar rooted at \(i\),
\begin{equation}
 a|_i=\sum_{X\ne\varnothing}a_{i,X}(w_X),\qquad
 a_{i,X}(0)=0,\qquad
 \sup_i\sum_X e^{\kappa\tau(\{i\}\cup X)+A|X|}
           \|a_{i,X}\|_{C^{1,\alpha};b}
       \le C_* b h .
 \label{lst-background}
\end{equation}
The derivatives in this norm are physical. Coefficients on
overlapping collars restrict one global assignment. All maxima,
first derivatives and first-gradient H\"older quotients are
inside the support sum. A fixed collar contains any coordinate
paths whose current length is at most \(L\); enlarging it changes
only fixed constants and possibly decreases \(\kappa\).
These are the complete-family inputs of the supported pre-gauge
background construction, not bounds merely on the summed field.

Let \(P_n(a)=K_n(a,\ldots,a)/n!\), fix \(a_0>0\), and set
\(c=a_0+1\). Choose \(J\) once with
\begin{equation}
 \frac{2c}{J}\le\kappa,\qquad A\ge4c .
 \label{lst-budgets}
\end{equation}

\begin{theorem}[Low-kernel spatial transport]
\label{lst-theorem}
There are \(C,L_0\), depending on
\(\sigma,J,C_*,a_0,\alpha\) and the fixed internal dimension,
but not on \(L,b\), volume or preceding refinement, such that
for \(L\ge L_0\) there are connected activities on final \(J\)-cells
with
\begin{equation}
 \sum_H B_H(w_H)=\sum_{n=2}^4P_n(a(w)),\qquad
 \sup_I\sum_{H\ni I}e^{a_0|H|}\|B_H\|_b
                      \le C\epsilon L^{-\alpha}.
 \label{lst-conclusion}
\end{equation}
The construction is linear in the total kernels.
Individual \(B_H\) need not be gauge invariant.
The scalar identity is vacuum-normalized and absolutely
convergent for compactly supported final fields.
\end{theorem}

The constant is fixed after \(J\) and before \(L\).
One can therefore set \(L=R^q\) with fixed \(R,J\) and
choose a finite \(q\) afterward. This statement does not
identify \(q\) successive quantum integrations with one
direct fluctuation integral.

\subsection{A valid quadratic difference factorization}
\label{lst-quadratic-section}

The matrix kernel need not be symmetric at each displacement,
even though Hessian reciprocity gives
\(K_2(r)^T=K_2(-r)\). We use a Laurent factorization that
does not require pointwise symmetry. Set
\[
 F(z)=\sum_r K_2(r)z^r,\qquad
 d_\rho^+(z)=z_\rho-1,\qquad
 d_\rho^-(z)=z_\rho^{-1}-1 .
\]
The exponential bound makes \(F\) holomorphic on a product
annulus. Equation~\eqref{lst-moments} says its jets through
degree three at \(z=(1,\ldots,1)\) vanish.

For a Laurent function define the ordered divided difference
\begin{equation}
 {\cal B}_jF(z)=
 \frac{F(1,\ldots,1,z_j,\ldots,z_4)
        -F(1,\ldots,1,1,z_{j+1},\ldots,z_4)}
      {z_j-1}.
 \label{lst-divided-difference}
\end{equation}
Its singularity at \(z_j=1\) is removable, and
\[
 F(z)-F(1)=\sum_j(z_j-1){\cal B}_jF(z).
\]
For a monomial \(z^r\), the quotient in coordinate \(j\)
has \(|r_j|\) monomials, each at absolute exponent length
at most \(|r|_1\), including negative exponents. Therefore,
for \(0<\sigma'<\sigma\), the combined absolute row Wiener
norm obeys
\begin{equation}
 \sup_u\sum_{j,v,r}e^{\sigma'|r|_1}
       |({\cal B}_jF)_{uv}(r)|
 \le\frac1{e(\sigma-\sigma')}\,
             \sup_u\sum_{v,r}e^{\sigma|r|_1}|F_{uv}(r)|.
 \label{lst-Wiener-quotient}
\end{equation}
Here the \(j\) sum is inside the row supremum. The bound
uses \(\sum_j|r_j|=|r|_1\) and
\(\sup_{t\ge0}t e^{-(\sigma-\sigma')t}
 =[e(\sigma-\sigma')]^{-1}\).
The local Taylor series also shows that
\({\cal B}_j\) lowers a vanishing order by at most one.

Apply the identity twice, using exponents
\(\sigma,3\sigma/4,\sigma/2\). Put
\[
 H_{\rho\tau}={\cal B}_\tau{\cal B}_\rho F,\qquad
 \widetilde H_{\rho\tau}(z)=-z_\rho H_{\rho\tau}(z).
\]
Since \(d_\rho^-(-z_\rho)=d_\rho^+\), this gives
\begin{equation}
 F(z)=\sum_{\rho,\tau}
       d_\rho^-(z)\widetilde H_{\rho\tau}(z)d_\tau^+(z),
 \qquad \widetilde H_{\rho\tau}(1)=0.
 \label{lst-factorization}
\end{equation}
In fact the jets of \(\widetilde H_{\rho\tau}\) through
degree one vanish. Its combined kernel bound is
\begin{equation}
 \sup_u\sum_{\rho,\tau,v,r}
     e^{(\sigma/2)|r|_1}
           |(\widetilde H_{\rho\tau})_{uv}(r)|
 \le \frac{16e^{\sigma/2}}{e^2\sigma^2}
                           \epsilon b^{-2}.
 \label{lst-H-bound}
\end{equation}
The factor \(e^{\sigma/2}\) pays the unit shift. No symmetry
of the individual \(\widetilde H_{\rho\tau}\) is asserted.

Let \(D_\rho a(x)=a(x+e_\rho)-a(x)\), the unscaled current
difference. Integration by parts transfers one difference
to each field. The zero sum in \eqref{lst-factorization}
then gives the exact identity
\begin{equation}
 P_2(a)=\frac12\sum_{x,r,\rho,\tau}
 [D_\rho a(x)]^T\widetilde H_{\rho\tau}(r)
            [D_\tau a(x+r)-D_\tau a(x)].
 \label{lst-quadratic-identity}
\end{equation}
The sign follows because \(d_\rho^-\) is the bilinear
adjoint of \(d_\rho^+\). Compactly supported fields give
the identity directly; the absolute majorants below
extend it to the substituted fields. No positivity or
second field derivative is used.

\subsection{Cubic and quartic cancellation}
\label{lst-higher-algebra}

Put
\begin{equation}
 \ell_r(x)=\sum_\rho r_\rho D_\rho a(x),\qquad
 R_r(x)=a(x+r)-a(x)-\ell_r(x).
 \label{lst-Taylor}
\end{equation}
In the anchored cubic expression
\[
 P_3(a)=\frac16\sum_{x,r,s}
                K_3(r,s)[a(x),a(x+r),a(x+s)],
\]
replace each shifted factor by \(a(x)+\ell+R\).
The zeroth moment cancels the all-anchor term. The
first moments cancel \(a,\ell,a\) and \(a,a,\ell\).
Six terms remain: the two with one \(R\) and no \(\ell\),
the \(\ell,\ell\) term, the two mixed \(\ell,R\) terms,
and the \(R,R\) term, always with unshifted first factor.

For the quartic expression subtract the all-anchor product
using its zeroth moment, then telescope the three shifted
factors. Each of three terms has one difference
\(a(x+r_j)-a(x)\) and three endpoint or anchor factors.
All cancellations are made in the total kernels before
triangle inequalities or final carrier assignment.

\subsection{Near-displacement coefficient bounds}
\label{lst-near-section}

Split these exact residual sums into \(m\le L\) and \(m>L\).
Neither restricted kernel is claimed to retain zero moments.
For the quadratic identity use \(m=|r|_1\).
In the near case all evaluations and coordinate paths lie
in a fixed final-unit collar. Coherence in
\eqref{lst-background} gives rooted coefficient-family bounds
\begin{align}
 |a|&\le Cbh,\qquad
 |\ell_r|\le Cbh^2|r|_1,\notag\\
 |a(x+r)-a(x)|&\le Cbh^2|r|_1,\qquad
 |R_r|\le Cbh^{2+\alpha}|r|_1^{1+\alpha}.
 \label{lst-near-fields}
\end{align}
For a negative coordinate traversal, the positive difference
is evaluated at the lower endpoint with the traversal sign.
Subtracting its value at the anchor and using the physical
gradient H\"older bound proves the last estimate.
Each inequality is for the weighted sum of coefficient
norms, with its suprema and H\"older quotients inside that
sum. It is not deduced from a bound on the scalar sum.

All near factors can be rooted at the anchor cell.
Multiplication glues rooted trees and support unions, so
their weights are submultiplicative. In the quadratic
identity,
\[
 |D_\rho a(x)|\le Cbh^2,\qquad
 |D_\tau a(x+r)-D_\tau a(x)|
                      \le Cbh^{2+\alpha}|r|_1^\alpha.
\]
Its product is therefore
\begin{equation}
 Cb^2h^{4+\alpha}|r|_1^\alpha.
 \label{lst-near-quadratic}
\end{equation}
Sum it with \eqref{lst-H-bound}. The six cubic products
are bounded by
\begin{equation}
 Cb^3h^{4+\alpha}(1+m)^6.
 \label{lst-near-cubic}
\end{equation}
The leading \(a,a,R\) term has that power of \(h\);
\(a,\ell,\ell\) has \(h^5\), and the others have still
higher powers. The quartic telescoping is bounded by
\begin{equation}
 Cb^4h^5(1+m).
 \label{lst-near-quartic}
\end{equation}
The weighted kernel sums cancel \(b^n\) and pay all
fixed displacement moments. Per original current anchor,
the three bounds are respectively
\(C\epsilon h^{4+\alpha}\),
\(C\epsilon h^{4+\alpha}\), and \(C\epsilon h^5\).

\subsection{Connected carriers and the anchor count}
\label{lst-carriers-section}

For each product coefficient keep its final-unit anchor
\(i_0\) and the endpoint roots actually used. In a near
term the complete collar permits the common root \(i_0\).
In a far term use the endpoint roots and any neighboring
roots needed for \(Da\). Their number is at most a fixed
\(n_0\), independent of \(L\) and the displacement.

Project the finite root set and, for each field-factor
occurrence, the terminal set consisting of its root
and variable support \(X\) to final \(J\)-cells.
For each projected terminal set choose a deterministic
short macro connecting graph supplied by the projection
lemma. Take the connected union of these macro graphs
as \(H\), including the anchor. The fine Steiner
lengths below bound those chosen macro graphs; the
carrier need not contain the full projection of an
arbitrary fine connecting tree. The projection bound gives
\begin{align}
 |H|&\le1+T,\notag\\
 T&\le\frac2J\left[
       \tau(\text{roots})+
       \sum_{\mathrm{occ}}\tau(\{\text{root}\}\cup X)\right]
         +4\left[n_0+\sum_{\mathrm{occ}}|X|\right].
 \label{lst-tree-bound}
\end{align}
Let \(j(x)\) be the final \(J\)-cell containing \(x/L\).
For \(z\in H\), its graph distance to \(j(x)\) is at
most \(T\). Hence
\begin{equation}
 e^{a_0|H|}1_{\{z\in H\}}
 \le e^{a_0}e^{-\operatorname{dist}(z,j(x))}e^{cT}.
 \label{lst-reroot}
\end{equation}
The budgets \eqref{lst-budgets} pay every attached
tree and support-cardinality factor. The fixed
\(e^{4cn_0}\) is absorbed into \(C\); near root-joining
trees have bounded length.

There are exactly \((JL)^4\) current anchor sites
per final \(J\)-cell, with the same centered-boundary
convention throughout. Therefore
\begin{equation}
 \sum_{x\in\mathbb Z^4}
 e^{-\operatorname{dist}(z,j(x))}
 \le (JL)^4 S_1,\qquad
 S_1=\sum_{y\in\mathbb Z^4}e^{-|y|_1}.
 \label{lst-anchor-count}
\end{equation}
The first internal-index count is the fixed \(v\).
There is no additional count of all output coordinates.
The near contributions are consequently
\(CJ^4\epsilon L^{-\alpha}\) for degrees two and
three, and \(CJ^4\epsilon L^{-1}\) for degree four.

\subsection{Far displacements and convergence}
\label{lst-far-section}

For \(m>L\), use the same exact residual formulas but
root each endpoint separately. The cubic and quartic
factors obey
\[
 |a|\le Cbh,\qquad
 |\ell_r|\le Cbh^2m,\qquad
 |R_r|\le Cbh(1+hm).
\]
Their products are bounded by
\(Cb^3h^3(1+m)^2\) and \(Cb^4h^4\).
For the quadratic expression both first differences
are bounded by \(Cbh^2\) at their own roots, so its
product is bounded by \(Cb^2h^4\).

The finite endpoint-root tree has
\[
 \tau(\text{roots})\le C_d(m/L+1).
\]
After \eqref{lst-reroot} this costs at most
\(C\exp(C'm/(JL))\). Use the common kernel exponent
\(\sigma_*=\sigma/2\), supplied by \eqref{lst-H-bound}
and also available for \(K_3,K_4\). Choose \(L_0\)
so that \(C'/(JL)\le\sigma_*/4\). One half of
\(\sigma_*\) pays the extra tree weight and the fixed
polynomial, leaving an exponentially small tail.
After \eqref{lst-anchor-count} the total far bound is
\begin{equation}
 CJ^4\epsilon(L+2)e^{-\sigma_*L/4}
                 \le C'\epsilon L^{-\alpha},
       \qquad L\ge L_0 .
 \label{lst-far-bound}
\end{equation}
All constants here are fixed before \(L\).

The absolute majorants justify regrouping by connected
\(H\) and give holomorphic coefficients on the same
open \(b\)-ball, with their bounded supremum norm.
For compactly supported \(w\), each nonzero field-factor
coefficient has variable support meeting that compact
set, since it vanishes at zero. Its final carrier
contains that support. The anchored bound therefore
implies convergence of the total sum, bounded by a
constant times the number of final \(J\)-cells meeting
the support. For that compact support \(K\), every
nonzero field coefficient also has
\(\tau(\{i\}\cup X)\ge\operatorname{dist}(i,K)\).
Equation~\eqref{lst-background} therefore gives
exponential spatial decay of \(a(w)\) and its first
differences. The original quadratic kernel sum is
absolutely convergent as well, so truncation and
exhaustion justify its integration by parts without
a boundary term. The near and far
bounds prove Theorem~\ref{lst-theorem}.

\subsection{Scope of the relocalization}
\label{lst-scope}

The result is a bounded connected localization of
the canceled low-degree \emph{total} functional.
It is stronger than extensive scalar power counting
and needs no invariant individual carriers. For a
residual analytic activity \(E-\beta S\), its
coefficientwise degree-\(\ge5\) part can be transported
by the previous high-jet estimate, while the total
degrees two through four use the present construction.
The two localizations may change the representation.

This section does not prove canonical finite-torus
wrapping control, an infinite-volume quantum measure,
equality of sequential and direct cutoff integrals,
boundedness of a complete quantum map on a quotient
of representations, or a coupling trajectory.
The classical scalar semigroup permits the choice
\(L=R^q\) for a gauge-invariant total, but removes
none of those separate interfaces.

% No preamble. Prefix aqb-. Infinite analytic interactions; classical descent only.
\section{A Banach quotient for analytic interactions}
\label{aqb-section}

Fix the cell size $J$, a bound $D$ on the number of scalar field
coordinates per cell, a cardinality exponent $a>0$, and a radius
$b>0$. Let $\mathcal C$ be the countable set of nonempty finite
connected cell sets. A declared carrier may exceed the variables
actually read by its coefficient. Let $H^\infty_b(H)$ denote the
bounded holomorphic functions on the open scalar-coordinate
radius-$b$ polydisc of $H$. Define
\begin{equation}
 \mathcal A_{a,b}=\left\{E=(E_H)_{H\in\mathcal C}:
 \|E\|_{a,b}:=\sup_i\sum_{H\ni i}e^{a|H|}
                       \|E_H\|_b<\infty\right\}.
 \label{aqb-activity-space}
\end{equation}
Individual coefficients need not have any symmetry. Constants
are allowed. All constants below may depend on the fixed $J,D,a$.

\subsection{Closed equality kernel and summable fields}
\label{aqb-quotient-section}

For a finitely supported field $A$, $\|A\|_\infty<b$, put
\begin{equation}
 \Phi_E(A)=\sum_H\bigl[E_H(A|_H)-E_H(0)\bigr].
 \label{aqb-total}
\end{equation}
If $K$ is its finite cell support, only carriers meeting $K$
contribute, and
\begin{equation}
 \sum_H|E_H(A|_H)-E_H(0)|\le2|K|\|E\|_{a,b}.
 \label{aqb-evaluation-bound}
\end{equation}
The series is normally convergent on every finite-coordinate
polydisc.

\begin{proposition}[Complete representation quotient]
\label{aqb-quotient-proposition}
The space $\mathcal A_{a,b}$ is Banach. The subspace
\begin{equation}
 \mathcal N_{a,b}=\{E:\Phi_E(A)=0
                  \text{ for every compact }A,\ \|A\|_\infty<b\}
 \label{aqb-null-space}
\end{equation}
is closed, and
\begin{equation}
 \mathcal Q_{a,b}=\mathcal A_{a,b}/\mathcal N_{a,b},
 \qquad
 \|[E]\|_{\mathcal Q}=\inf_{F\in\mathcal N_{a,b}}
                                      \|E+F\|_{a,b}
 \label{aqb-quotient-norm}
\end{equation}
is a Banach space. Its invariant subspace, defined by translation,
hypercubic, global-adjoint and local-gauge invariance of the
normalized total, is closed.
\end{proposition}

\begin{proof}
A norm-Cauchy sequence is Cauchy in each $H^\infty_b(H)$:
choose one root in $H$ in \eqref{aqb-activity-space}. Component
limits exist. Fatou over carriers, followed by the root supremum,
proves both the norm bound and norm convergence. The null space
is the intersection of the kernels of the bounded evaluations
\eqref{aqb-evaluation-bound}. The usual quotient-completeness
argument uses a summably Cauchy subsequence and representatives
with summable successive differences; it requires no minimizing
representative. Each total symmetry is an equality of bounded
evaluations, for all compact configurations and transforms in
their common domain. Its preimage is closed and contains
$\mathcal N_{a,b}$. Reality specifies a closed real form.
\end{proof}

The quotient removes all carrierwise constants and all other
representations of the zero normalized total. It does not select
a physical vacuum counterterm; a vacuum sector, if needed, stays
separate. Equality of germs on every finite-coordinate subspace
is equivalent to \eqref{aqb-null-space} by analytic uniqueness.

Use the cell norm $\|A\|_1=\sum_i\|A|_i\|_\infty$.
For $\|A\|_\infty,\|B\|_\infty\le\rho<b$, directional
Cauchy on the convex segment gives
\begin{align}
 |E_H(A|_H)-E_H(B|_H)|
 &\le\frac{\|E_H\|_b}{b-\rho}
                          \|(A-B)|_H\|_\infty\notag\\
 &\le\frac{\|E_H\|_b}{b-\rho}
                          \sum_{i\in H}\|A_i-B_i\|_\infty.
 \label{aqb-local-Lipschitz}
\end{align}
Use a strictly smaller directional circle first, then take its
radius to the displayed margin. Summing at each cell proves
\begin{equation}
 |\Phi_E(A)-\Phi_E(B)|
       \le\frac{\|E\|_{a,b}}{b-\rho}\|A-B\|_1.
 \label{aqb-summable-Lipschitz}
\end{equation}
Thus \eqref{aqb-total} extends by an absolutely convergent series
to summable fields in every strict interior ball. Finite field
truncations converge in \eqref{aqb-summable-Lipschitz}. A null
family remains zero on these fields. The extension therefore
depends only on the quotient class, and the estimate holds with
the quotient norm after taking the infimum over representatives.

\subsection{A continuous total Maxwell coefficient}
\label{aqb-beta-section}

For each fixed order $n\le4$, scalar Cauchy and the exact finite
monomial count give a summable total derivative kernel. Indeed
the ordered derivative sum on a carrier is at most
$(D|H|+n)^n b^{-n}\|E_H\|_b$. Connectedness bounds every
slot displacement by $c_dJ|H|$. Spending half the carrier
exponent on spatial distance and half on the fixed polynomial
gives, for sufficiently small fixed $\sigma=\sigma(a,J)>0$,
\begin{equation}
 \sup_{i_1}\sum_{i_2,\ldots,i_n}
 e^{\sigma\sum_{j\ge2}|x_j-x_1|}
 |K_n(E;i_1,\ldots,i_n)|
       \le C_n(a,D,J)b^{-n}\|E\|_{a,b}.
 \label{aqb-kernel-bound}
\end{equation}
All scalar and color slots are included. Fixed moments follow
with a further exponent loss. This is a fixed-order estimate,
not a dimension-free absolute Taylor estimate at all orders.

Each kernel component is a derivative on a finite-coordinate
subspace and hence depends only on $[E]$. Taking a representative
infimum gives \eqref{aqb-kernel-bound} with $\|[E]\|_{\mathcal Q}$.
For a simple gauge group and an invariant total, define in the
trace-normalized color component
\begin{equation}
 \beta([E])=\frac12\partial_{p_\nu}^2
                       \Gamma^E_{2,\mu\mu}(0),
 \qquad \mu\ne\nu.
 \label{aqb-beta-definition}
\end{equation}
Use positive-bond midpoints $\xi(x,\mu)=x+e_\mu/2$.
The total Ward algebra of Theorem~\ref{tkw-theorem} applies
directly to these infinite kernels. It identifies
\eqref{aqb-beta-definition} as the Maxwell coefficient and makes
the distinct pair immaterial. No subsequence is needed once
the infinite interaction is specified. For real interactions
this coefficient is real.

It is a continuous linear functional on the invariant quotient:
\begin{equation}
 |\beta([E])|\le C_\beta(a,D,J)b^{-2}
                                      \|[E]\|_{\mathcal Q}.
 \label{aqb-beta-bound}
\end{equation}
For example, a second derivative component is bounded by
$2b^{-2}\|E_H\|_b$, there are at most $D|H|$ second slots,
and their squared midpoint separation is at most
$(c_dJ|H|)^2$. Thus, up to a fixed color-normalization factor,
one can take
\begin{equation}
 C_\beta=D(c_dJ)^2\sup_{m\ge1}m^3e^{-am}.
 \label{aqb-beta-constant}
\end{equation}
The factor $b^{-2}$ is essential. For a semisimple product use
the finite vector of coefficients and one reference per simple
ideal.

\subsection{Full-reference extraction without a small-radius loss}
\label{aqb-extraction-section}

Let $[S]$ be an invariant reference with $\beta([S])=1$.
Assume a representation on one fixed radius $\rho$ satisfies
$\|S\|_{a,\rho}\le M_S$, uniformly in preceding refinement,
and its total first derivative vanishes. This is the uniform
fixed-radius classical support input; normalization of the
Hessian alone does not imply it.

Define the homogeneous phase projections coefficientwise by
\begin{equation}
 (P_nE_H)(A)=\frac1{2\pi}\int_0^{2\pi}
                     e^{-int}E_H(e^{it}A)\,dt.
 \label{aqb-phase-projection}
\end{equation}
They retain $H$ and have norm at most one. Replace $S_H$ by
$S_H^{\ge2}=\sum_{n\ge2}P_nS_H$. The normalized total has
the same class, since constants are removed and the total
linear sector is zero. Homogeneity yields, for $b\le\rho/2$,
\begin{equation}
 \|S^{\ge2}\|_{a,b}
 \le\frac{(b/\rho)^2}{1-b/\rho}M_S
 \le C_Sb^2,
 \qquad C_S=\frac{2M_S}{\rho^2}.
 \label{aqb-reference-bound}
\end{equation}
No vanishing of individual linear coefficients or individual
gauge invariance was assumed. A joint subsequential limit
preserves this bound by Fatou.

The full-reference projection
\begin{equation}
 \mathcal P_S[E]=[E]-\beta([E])[S]
 \label{aqb-extraction}
\end{equation}
satisfies $\beta\mathcal P_S=0$, $\mathcal P_S^2=\mathcal P_S$,
and
\begin{equation}
 \|\mathcal P_S\|\le C_{\rm ext},
 \qquad C_{\rm ext}=1+C_\beta C_S.
 \label{aqb-extraction-bound}
\end{equation}
The constant is independent of small $b$. Apply
\eqref{aqb-beta-bound} and \eqref{aqb-reference-bound} to
representatives approaching the quotient infimum. The whole
invariant reference is subtracted, so its nonlinear Ward
identities remain exact.

\subsection{Classical substitution on the quotient}
\label{aqb-descent-section}

Fix $L$ and an actual infinite pre-gauge background map $a_L$.
Assume its coherent whole-fiber support family satisfies
\begin{equation}
 a_L|_i=\sum_{X\ne\varnothing}a_{i,X}(w_X),\qquad
 a_{i,X}(0)=0,\qquad
 \sup_i\sum_X e^{\kappa\tau(\{i\}\cup X)+A|X|}
                \|a_{i,X}\|_b\le Cb/L,
 \label{aqb-background-family}
\end{equation}
and $\|a_L(w)\|_\infty\le\rho_0<b$. The maximum over the
whole current output fiber is inside the support sum, and
overlapping collars use one coefficient assignment. For compact
$w$ supported in $K$, a nonzero term has $X\cap K\ne\varnothing$;
hence
\begin{equation}
 \|a_L(w)|_i\|\le(Cb/L)e^{-\kappa\operatorname{dist}(i,K)}.
 \label{aqb-background-summability}
\end{equation}
At fixed $L$ there are finitely many current coordinates per
fiber, so $a_L(w)$ is summable. Consequently
\begin{equation}
 (\mathcal T_L\Phi_E)(w)=\Phi_E(a_L(w))
 \label{aqb-scalar-substitution}
\end{equation}
depends only on $[E]$, by \eqref{aqb-summable-Lipschitz}.
This is algebraic classical descent.

For bounded descent into an activity quotient, suppose a
connected strict representation $\mathscr L_LE$ of
\eqref{aqb-scalar-substitution} obeys
\begin{equation}
 \|\mathscr L_LE\|_{a_{\rm out},b}
                         \le B_L\|E\|_{a,b}.
 \label{aqb-representation-input}
\end{equation}
Then $\mathcal T_L[E]=[\mathscr L_LE]$ is well defined and has
norm at most $B_L$. Distinct valid output assignments have the
same normalized total; equivalent inputs give equal output
totals by summable-field continuity. The class map is linear
because its scalar total is linear in $E$, even without a linear
choice of representatives. Taking near-minimal input
representatives proves the norm bound. Scalar amplitude alone
does not prove \eqref{aqb-representation-input}.

The compactness construction of Section~\ref{aic-section}
provides simultaneous infinite interactions and coherent maps
under its fixed-parameter finite-regulator hypotheses. Its
summable-field composition limit transfers finite classical
identities only on one jointly chosen system, with common strict
domain margins. Unrelated subsequences cannot be composed.
Preserving invariant totals needs more than an abstract orbit
relation. The connecting infinite gauge must be one to which
the input total's gauge identity applies. A sufficient condition
is approximation by allowed compact gauges with transformed
fields converging in $\ell^1$ inside a common strict margin;
then \eqref{aqb-summable-Lipschitz} passes the identity. A joint
finite invariant approximation of a particular interaction and
the maps is another sufficient condition for that interaction.
For an action on the entire invariant quotient, this requirement
must hold for every input class. A joint approximation of one
input, or unrelated finite subsequences, does not establish a
universal invariant-quotient action.

\subsection{Conditional classical contraction after extraction}
\label{aqb-contraction-section}

The following statement keeps the infinite-family and scale
requirements explicit. Fix $J,a$ and the primitive block ratio.
Use one joint infinite system with normalized references $S_M$
and $S_{ML}$ satisfying \eqref{aqb-reference-bound} uniformly.
Assume the coherent physical $C^{1,\alpha}$ families of
Proposition~\ref{pgh-support-proposition}, with primitive
constants independent of $L,M$, and impose the fixed support
reserves of Theorem~\ref{lst-theorem} after choosing $J$.
Assume also the core/tail budgets of the fifth-order background
transport theorem of Notebook28, with unchanged $a,b$.
Its constant $C_{\rm high}$ is independent of $L$ once these
reserves are fixed. Increasing $L$ weakens its compression
requirement; the same fixed $J$ and logarithmic tail cutoff
remain admissible. Choose $b$ after the final finite $L$, so
$b\le r_A/L_*$ with $L_*=L/(8C_A)$ and all other classical
margins hold. Scalar substitution and preservation of the
invariant domain satisfy the additional gauge-identity requirement
just specified. An operator assertion on the whole invariant
quotient assumes it for every class; a joint finite approximation
may instead verify the following pointwise estimate on its
specified collection of classes.

\begin{theorem}[Classical quotient estimate]
\label{aqb-classical-theorem}
Under these hypotheses, in dimension four,
\begin{equation}
 \|\mathcal T_L\mathcal P_{S_M}[E]\|_{\mathcal Q}
 \le C_{{\rm ext,in}}
       \bigl[C_{\rm low}(J)L^{-\alpha}
                    +6C_{\rm high}L^{-1}\bigr]
                     \|[E]\|_{\mathcal Q}.
 \label{aqb-classical-bound}
\end{equation}
If a normalized output reference is used, then
\begin{equation}
 \|\mathcal P_{S_{ML}}\mathcal T_L\mathcal P_{S_M}\|
 \le C_{{\rm ext,out}}C_{{\rm ext,in}}
       \bigl[C_{\rm low}(J)L^{-\alpha}
                    +6C_{\rm high}L^{-1}\bigr].
 \label{aqb-projected-contraction}
\end{equation}
The constants are fixed after $J$ and before $L,b$. A sufficiently
large finite allowed $L=R^q$, followed by a sufficiently small
common positive $b$, makes the last bound less than one.
\end{theorem}

\begin{proof}
Put $[F]=\mathcal P_{S_M}[E]$ and choose any representative.
The phase projections \eqref{aqb-phase-projection} retain carriers,
have norm at most one, and give a coefficientwise fifth-order
remainder $R_5=I-\sum_{n=0}^4P_n$ with norm at most six.
These projections also descend to the unrestricted quotient:
apply the phase integral to a zero normalized analytic total.
They need not preserve gauge invariance separately.

The total linear term of $F$ is zero by global adjoint invariance,
its vacuum is removed, and $\beta(F)=0$. The total Ward algebra
therefore gives the low-kernel cancellations in
Theorem~\ref{lst-theorem}. The kernel bound
\eqref{aqb-kernel-bound} supplies that theorem's $\epsilon$ with
at most a fixed $C_{\rm ker}(J)\|F\|_{a,b}$. Its new connected
representation has norm at most
$C_{\rm low}(J)L^{-\alpha}\|F\|_{a,b}$.
The coefficientwise $R_5$ part has a connected representation
of norm at most $6C_{\rm high}L^{-1}\|F\|_{a,b}$ by the
fifth-order transport theorem. The zero total linear term stays
zero after substitution by \eqref{aqb-summable-Lipschitz}; it
need not vanish on each carrier.

All scalar decompositions are legitimate on summable inputs:
the phase projections and $R_5$ lie in the anchored space,
Section~\ref{aqb-quotient-section} gives their continuous
extensions, and the new output families converge normally on
compact fields. Adding the two representations and taking the
infimum over representatives of $[F]$ gives the bracket in
\eqref{aqb-classical-bound}. Apply
\eqref{aqb-extraction-bound} at the input, and if desired at
the output, to obtain both displays.
\end{proof}

The output projection is explicit: the theorem does not claim
that classical substitution itself preserves $\beta=0$.
The scalar semigroup of Theorem~\ref{csg-theorem}, passed on
a joint system as in Section~\ref{aic-section}, permits the
finite choice $L=R^q$ after $R,J$ are fixed. The bound is on a
class admitting a newly justified localization, not on an
arbitrarily fixed sequential carrier assignment.

\subsection{Quantum descent is a separate obligation}
\label{aqb-scope}

No quantum map on $\mathcal Q_{a,b}$ has been defined here.
Equality of infinite normalized totals on compact fields does
not by itself identify finite-regulator exponentials, their
sharp restricted integrals, chosen vacuum normalizations, or
their limits. Thus this quotient does not yet permit arbitrary
equivalent representatives to replace the actual quantum output
inside its next finite integration. Quantum descent and its
continuity and normalization compatibility, the full amplitude
recursion, and the coupling trajectory remain separate.

\section{The interface with a renormalized quantum induction}
\label{gex-next-section}

The low-degree classification is now tied to the actual local
symmetries and a normalized classical reference. The exact
coupling rebase is
\[
 g_{\mathrm{next}}^{-2}=g^{-2}-\beta,\qquad
 E_{\mathrm{ren}}=E-E(0)-\beta S_{MR}.
\]
This is an identity of the density after extraction. A sign
or trajectory estimate for the generated \(\beta\) is still
needed; the identity does not supply one.

There is a concrete route to the spatial part of a contraction.
Fix the elementary macrostep \(R\) and its sufficient auxiliary
grouping \(J\). The spatial remainder estimate has constants
depending on these fixed data. Select a finite \(q\) afterward,
so that its endpoint ratio \(L=R^q\) makes
\(C(J)L^{-\alpha}\) small. The classical semigroup identifies
the scalar inherited background with the direct endpoint
background. Its new connected localization can then be
estimated directly. The high-degree piece has the separate
inverse-\(L\) estimate of Notebook 28, with its analytic radius
chosen after this finite \(L\).

The analytic compactness and quotient constructions make
this classical change of representation precise. They supply
joint limiting interactions and maps, a complete quotient,
and a bounded full-reference projection. The final classical
estimate is a contraction once the finite endpoint ratio
is sufficiently large and its common radius is chosen.

To turn this into a quantum induction, four interfaces must
be established. First, the limiting interaction space must
be matched to the regulators used by the actual quantum update,
including winding terms and its vacuum normalization.
Second, prove that the quantum update descends continuously
to the quotient; classical descent alone does not identify
the sharp finite-regulator integrals for equivalent inputs.
Third, the \(q\) successive
quantum steps must retain all intermediate generated terms,
cutoffs and normalizations; the classical semigroup does not
replace them by a single fluctuation integral. Fourth, the
resulting recurrence must control both the renormalized
activity and the actual coupling trajectory uniformly in
the number of blocks of steps.

The present results supply the nonlinear marginal identity,
the classical spatial gain and its representation space
under stated hypotheses. The larger existence proof
still needs the mixed and large-field histories, sourced
observable estimates, continuum nontriviality, Euclidean
symmetry and reconstruction already recorded in Notebook 01.
The mass gap remains a subsequent spectral assertion.

\begin{thebibliography}{9}
\bibitem{priorgex}
\emph{Toward a four-dimensional Yang--Mills construction},
research Notebooks 07, 15--16 and 25--28 in this workspace.
The flat classical normalization, ordinary minimum theorem,
strict classical support families and common-domain quantum
macrostep are the principal prior inputs.
\bibitem{balaban109gex}
T.~Ba\l aban,
\emph{Renormalization group approach to lattice gauge field
theories. I. Generation of effective actions in a small field
approximation and a coupling constant renormalization in four
dimensions},
Commun. Math. Phys. \textbf{109} (1987), 249--301.
\href{https://doi.org/10.1007/BF01215223}{Publisher record};
\href{https://scispace.com/pdf/renormalization-group-approach-to-lattice-gauge-field-3ygai0vxt1.pdf}{Original article copy}.
The averaging symmetries and total Ward/subtraction calculations
are checked against this primary source. Its stronger
coefficientwise orbit hypotheses are not assumed here.
\bibitem{balaban116gex}
T.~Ba\l aban,
\emph{Renormalization group approach to lattice gauge field
theories. II. Cluster expansions},
Commun. Math. Phys. \textbf{116} (1988), 1--22.
\href{https://doi.org/10.1007/BF01239022}{Publisher record}.
\bibitem{sourcegex}
E.~Sussman and collaborators,
\href{https://github.com/ethanwsussman/Yang--Mills-LLM-experiment}
{Yang--Mills-LLM-experiment},
deposited raw accounts at commit
\texttt{e392ef45013a581664bab76b90867170b30fe10e}.
\end{thebibliography}
\end{document}
